\documentclass[
    11pt,
    a4paper,
    oneside,
    open=any,
    headings=normal
]{scrbook}

% ==================================================
% Energietechnik – modularer Praxiskurs
% Kompilierung: LuaLaTeX (z. B. zweimal ausführen)
% ==================================================

\usepackage{fontspec}
\usepackage[ngerman]{babel}
\usepackage{microtype}
\usepackage{geometry}
\geometry{left=28mm,right=25mm,top=25mm,bottom=27mm,headheight=15pt}

\usepackage{amsmath,amssymb,mathtools}
\usepackage{siunitx}
\sisetup{
  locale = DE,
  per-mode = symbol,
  separate-uncertainty = true
}

\usepackage{graphicx}
\usepackage{tikz}
\usetikzlibrary{
  arrows.meta,
  calc,
  circuits.ee.IEC,
  positioning,
  shapes.geometric
}
\usepackage[european]{circuitikz}

\usepackage{booktabs}
\usepackage{tabularx}
\usepackage{enumitem}
\usepackage{csquotes}
\usepackage{xcolor}
\usepackage{fancyhdr}
\usepackage[most]{tcolorbox}
\usepackage{hyperref}
\usepackage[nameinlink,noabbrev]{cleveref}

\definecolor{NetzBlau}{HTML}{174A70}
\definecolor{WarnRot}{HTML}{A52A2A}
\definecolor{HinweisGrau}{HTML}{F0F2F4}
\definecolor{SicherheitsGelb}{HTML}{FFF4CC}

\definecolor{Seitenfarbe}{HTML}{2f4f67}

\hypersetup{
  colorlinks=true,
  linkcolor=NetzBlau,
  citecolor=NetzBlau,
  urlcolor=NetzBlau,
  pdftitle={Energietechnik – Netze, Schutz, Prüfungen und Betriebssicherheit},
  pdfsubject={Technisches Lehrskript}
}

\pagestyle{fancy}
\fancyhf{}
\fancyhead[L]{\nouppercase{\leftmark}}
\fancyhead[R]{Energietechnik – Praxiskurs}
\fancyfoot[C]{\thepage}
\setlength{\headheight}{15pt}

% --------------------------------------------------
% Einheitliche Lernbausteine
% --------------------------------------------------
\newcommand{\phasor}[1]{\underline{#1}}
\newcommand{\jj}{\mathrm{j}}
\newcommand{\dd}{\,\mathrm{d}}
\newcommand{\Ohm}{\,\si{\ohm}}
\newcommand{\kV}{\,\si{\kilo\volt}}
\newcommand{\A}{\,\si{\ampere}}
\newcommand{\kA}{\,\si{\kilo\ampere}}

\newenvironment{lernziel}
{}{}
  % {\par\medskip\noindent\colorbox{HinweisGrau}{%
  %   \begin{minipage}{0.94\linewidth}\textbf{Lernziele}\par\smallskip}}
  % {\end{minipage}\par\medskip}

\newenvironment{merke}
{}{}
  % {\par\medskip\noindent\colorbox{HinweisGrau}{%
  %   \begin{minipage}{0.94\linewidth}\textbf{Merke}\par\smallskip}}
  % {\end{minipage}\par\medskip}

\newtcolorbox{sicherheitshinweis}{
  title=\textbf{\color{black} Sicherheitshinweis},
  colback=SicherheitsGelb,
  colframe=SicherheitsGelb,
  boxrule=0pt,
  arc=0pt,
  left=2mm,
  right=2mm,
  top=2mm,
  bottom=2mm,
  before skip=\medskipamount,
  after skip=\medskipamount,
  fontupper=\normalsize
}

% \newenvironment{sicherheitshinweis}
% {
%   \par\medskip\noindent\colorbox{SicherheitsGelb}
%   {Rechenbeispiel}
% }{

% }

  % {\par\medskip\noindent\colorbox{SicherheitsGelb}{%
  %   \begin{minipage}{0.94\linewidth}\textbf{Sicherheitshinweis}\par\smallskip}}
  % {\end{minipage}\par\medskip}

\newenvironment{rechenbeispiel}
  {\par\medskip\noindent\textbf{Rechenbeispiel}\par\smallskip}
  {\par\medskip}

\newenvironment{aufgaben}
  {\par\medskip\noindent\textbf{Übungsaufgaben}\begin{enumerate}[label=\arabic*.]}
  {\end{enumerate}\medskip}

% --------------------------------------------------
% Metadaten
% --------------------------------------------------
\title{\Huge Energietechnik\\[0.5em]
  \Large Netze, Schutz, Prüfungen und Betriebssicherheit}
\subtitle{Vertiefter Praxiskurs mit Berechnungen und Fallstudien}
\author{Technisches Lehrskript}
\date{\today}

\begin{document}
\frontmatter
\maketitle

%\pagecolor{Seitenfarbe}
%\color{lightgray}

\chapter*{Ziel und Arbeitsweise}
Dieses Lehrskript verbindet die physikalischen Grundlagen der elektrischen
Energietechnik mit Netzberechnung, Betriebsmitteln, Schutztechnik, Prüfungen
und sicherem Anlagenbetrieb. Es ist als modular wachsendes Arbeitsdokument
angelegt.

Jedes Fachkapitel soll nach Möglichkeit enthalten:
\begin{itemize}
  \item physikalisches Modell und Voraussetzungen,
  \item Ersatzschaltbild und Definition der Bezugsgrößen,
  \item nachvollziehbare Herleitung,
  \item Zahlenbeispiel mit Einheitenprüfung,
  \item typische Denk- und Rechenfehler,
  \item Übungsaufgaben mit Musterlösungen,
  \item Hinweise auf relevante Normen und deren Geltungsbereich.
\end{itemize}

\begin{sicherheitshinweis}
Dieses Skript ist eine Lernunterlage. Es ersetzt weder die jeweils gültigen
Normen und betrieblichen Anweisungen noch Qualifikation, Gefährdungsbeurteilung,
Anlagenkenntnis oder eine konkrete Freigabe. Schalthandlungen und Arbeiten an
elektrischen Anlagen dürfen ausschließlich durch dafür qualifizierte und
beauftragte Personen nach den gültigen Verfahren erfolgen.
\end{sicherheitshinweis}

\tableofcontents

\mainmatter

\part{Grundlagen der elektrischen Energietechnik}
\chapter{Drehstrom und komplexe Wechselstromrechnung}
\label{chap:drehstrom}

\begin{lernziel}
Nach diesem Kapitel kannst du:
\begin{itemize}
  \item Effektivwert, Scheitelwert, Frequenz und Phasenwinkel unterscheiden,
  \item sinusförmige Spannungen und Ströme als Zeigergrößen darstellen,
  \item Impedanzen von Widerstand, Induktivität und Kapazität bilden,
  \item Wirk-, Blind- und Scheinleistung voneinander unterscheiden,
  \item Leiter- und Stranggrößen im symmetrischen Drehstromsystem umrechnen.
\end{itemize}
\end{lernziel}

\section{Warum wir mit Wechselstromrechnung beginnen}

Netze der öffentlichen Energieversorgung arbeiten überwiegend mit
sinusförmigem Drehstrom. Auch wenn reale Netze Oberschwingungen, Unsymmetrien
und transiente Vorgänge enthalten, ist das symmetrische sinusförmige
Grundmodell der Ausgangspunkt für viele Netzberechnungen.

Die zentrale Vereinfachung lautet: Im eingeschwungenen sinusförmigen Zustand
mit einer festen Frequenz können Spannungen und Ströme durch komplexe
Zeigergrößen und Bauteile durch komplexe Impedanzen beschrieben werden.
Differentialgleichungen werden damit für viele stationäre Berechnungen zu
algebraischen Gleichungen.

\section{Sinusförmige Größen}

Eine sinusförmige Spannung kann als
\begin{equation}
  u(t)=\hat U\cos(\omega t+\varphi_u)
  \label{eq:sinusspannung}
\end{equation}
geschrieben werden. Dabei ist:
\begin{description}[style=nextline]
  \item[$u(t)$] Momentanwert der Spannung,
  \item[$\hat U$] Scheitelwert,
  \item[$\omega$] Kreisfrequenz in \si{\radian\per\second},
  \item[$t$] Zeit,
  \item[$\varphi_u$] Phasenwinkel zum gewählten Zeitbezug.
\end{description}

Frequenz und Kreisfrequenz hängen zusammen:
\begin{equation}
  \omega=2\pi f.
\end{equation}
Für eine sinusförmige Größe gilt zwischen Scheitelwert und Effektivwert:
\begin{equation}
  U=\frac{\hat U}{\sqrt{2}},
  \qquad
  I=\frac{\hat I}{\sqrt{2}}.
\end{equation}
Der Effektivwert ist derjenige Gleichspannungs- beziehungsweise Gleichstromwert,
der an einem ohmschen Widerstand dieselbe mittlere Wärmeleistung erzeugen würde.

\begin{merke}
In Netzangaben wie \SI{400}{\volt} oder \SI{20}{\kilo\volt} ist im üblichen
Sinusmodell der Effektivwert gemeint, sofern nichts anderes angegeben ist.
Scheitelwerte dürfen nicht unbemerkt mit Effektivwerten vermischt werden.
\end{merke}

\section{Komplexe Zeigergrößen}

Wir verwenden im Folgenden die Kosinuskonvention. Eine sinusförmige Größe
\[
x(t)=\hat X\cos(\omega t+\varphi_x)
\]
wird durch den Effektivwertzeiger
\begin{equation}
  \phasor{X}=X e^{\jj\varphi_x}
  =X(\cos\varphi_x+\jj\sin\varphi_x)
\end{equation}
repräsentiert. Der Zeitverlauf lässt sich daraus mit
\[
x(t)=\sqrt{2}\,\operatorname{Re}
\left\{\phasor{X}e^{\jj\omega t}\right\}
\]
rekonstruieren.

Die komplexe Schreibweise trennt Betrag und Phasenlage. Das ist entscheidend,
weil an Induktivitäten und Kapazitäten Strom und Spannung nicht phasengleich
sind.

\subsection{Impedanz}

Im sinusförmigen stationären Zustand ist die Impedanz das Verhältnis aus
Spannungs- und Stromzeiger:
\begin{equation}
  Z=\frac{\phasor{U}}{\phasor{I}}=R+\jj X.
  \label{eq:impedanz}
\end{equation}
Hierbei bezeichnet \(R\) den Wirkwiderstand und \(X\) die Reaktanz.
Der Betrag ist
\[
|Z|=\sqrt{R^2+X^2},
\]
und der Impedanzwinkel beträgt
\[
\varphi_Z=\arg(Z)=\arctan\!\left(\frac{X}{R}\right),
\]
wobei für eine robuste Quadrantenwahl rechnerisch die Funktion
\(\operatorname{atan2}(X,R)\) zu verwenden ist.

Für die drei idealisierten Grundelemente gilt:
\begin{align}
  Z_R &= R,\\
  Z_L &= \jj\omega L,\\
  Z_C &= \frac{1}{\jj\omega C}=-\frac{\jj}{\omega C}.
\end{align}
Die induktive Reaktanz \(X_L=\omega L\) wächst mit der Frequenz.
Der Betrag der kapazitiven Reaktanz \( |X_C|=1/(\omega C) \) sinkt mit
wachsender Frequenz.

\section{Wirk-, Blind- und Scheinleistung}

Für Effektivwertzeiger verwenden wir die komplexe Leistung
\begin{equation}
  \underline S=\phasor U\,\phasor I^{*}=P+\jj Q,
  \label{eq:komplexeleistung}
\end{equation}
wobei \(^{*}\) die komplexe Konjugation bezeichnet. Daraus folgen:
\begin{align}
  P &= UI\cos\varphi && \text{Wirkleistung},\\
  Q &= UI\sin\varphi && \text{Blindleistung},\\
  S &= UI=\sqrt{P^2+Q^2} && \text{Scheinleistung}.
\end{align}
Die Vorzeichenkonvention für \(Q\) muss bei jeder konkreten Aufgabe klar sein.
In der üblichen Verbraucherzählpfeilkonvention ist induktive Blindleistung
einer Last positiv und kapazitive Blindleistung negativ.

\section{Symmetrisches Drehstromsystem}

Ein symmetrisches Drehstromsystem besitzt drei sinusförmige Strangspannungen
gleicher Frequenz und gleichen Effektivwerts, die zeitlich um jeweils
\SI{120}{\degree} verschoben sind. Bei positiver Phasenfolge \(L1,L2,L3\)
können die Strangspannungszeiger gegen den Sternpunkt geschrieben werden als
\begin{align}
  \phasor U_{1N}&=U_\mathrm{Str}\angle 0^\circ,\\
  \phasor U_{2N}&=U_\mathrm{Str}\angle -120^\circ,\\
  \phasor U_{3N}&=U_\mathrm{Str}\angle +120^\circ.
\end{align}
Ihre vektorielle Summe ist null:
\[
\phasor U_{1N}+\phasor U_{2N}+\phasor U_{3N}=0.
\]

Für ein symmetrisches System gilt zwischen Leiter-Leiter-Spannung und
Strangspannung:
\begin{equation}
  U_\mathrm{LL}=\sqrt{3}\,U_\mathrm{Str}.
  \label{eq:leiter_strang_spannung}
\end{equation}
Daher entsprechen bei einem üblichen \SI{400}{\volt}-Netz ungefähr
\SI{230}{\volt} zwischen Außenleiter und Neutralleiter. Bei einem
\SI{20}{\kilo\volt}-Netz beträgt die Strangspannung gegen den geerdeten
Sternpunkt im idealisierten Betrieb etwa \SI{11.55}{\kilo\volt}.

\begin{merke}
Die Beziehung \(U_\mathrm{LL}=\sqrt{3}U_\mathrm{Str}\) gilt für das
symmetrische Drehstromsystem. Bei Erdschluss und Sternpunktverschiebung darf
man die tatsächlichen Leiter-Erde-Spannungen nicht unkritisch aus dieser
Normalbetriebsbeziehung ableiten.
\end{merke}

\section{Drehstromleistung}

Für ein symmetrisches Drehstromsystem gilt:
\begin{align}
  P &= \sqrt{3}\,U_\mathrm{LL} I_\mathrm{L}\cos\varphi,\\
  Q &= \sqrt{3}\,U_\mathrm{LL} I_\mathrm{L}\sin\varphi,\\
  S &= \sqrt{3}\,U_\mathrm{LL} I_\mathrm{L}.
\end{align}
Dabei sind \(U_\mathrm{LL}\) die Leiter-Leiter-Spannung und \(I_\mathrm{L}\)
der Leiterstrom. Diese Formeln beziehen sich auf einen symmetrischen Zustand.

\begin{rechenbeispiel}
Ein symmetrischer Verbraucher wird an einem \SI{400}{\volt}-Netz betrieben.
Der Leiterstrom beträgt \SI{32}{\ampere}, der Leistungsfaktor ist
\(\cos\varphi=0{,}90\). Die aufgenommene Wirkleistung ist
\begin{align*}
  P &= \sqrt{3}\cdot \SI{400}{\volt}\cdot
       \SI{32}{\ampere}\cdot 0{,}90\\
    &\approx \SI{19.95}{\kilo\watt}.
\end{align*}
Die Scheinleistung beträgt
\[
S=\sqrt{3}\cdot\SI{400}{\volt}\cdot\SI{32}{\ampere}
 \approx \SI{22.17}{\kilo\volt\ampere}.
\]
Die Blindleistung ergibt sich aus \(Q=\sqrt{S^2-P^2}\), sofern der
induktive beziehungsweise kapazitive Charakter separat bekannt ist:
\[
Q\approx \SI{9.65}{\kilo var}
\]
bei induktivem Verbraucher und der hier angenommenen Vorzeichenkonvention.
\end{rechenbeispiel}

\section{Typische Fehler}

\begin{itemize}
  \item Leiter-Leiter-Spannung und Strangspannung verwechseln.
  \item Effektivwert und Scheitelwert vermischen.
  \item Bei komplexen Größen nur Beträge rechnen und Phasenwinkel ignorieren.
  \item Blindleistung mit Wirkleistung addieren, als wären beide dieselbe Art
        von Größe.
  \item Die Drehstromleistungsformeln auf unsymmetrische Betriebszustände
        anwenden, ohne die Voraussetzungen zu prüfen.
\end{itemize}

\begin{aufgaben}
  \item Eine sinusförmige Spannung hat den Scheitelwert
        \(\hat U=\SI{325}{\volt}\). Bestimme Effektivwert und Frequenz,
        wenn die Periodendauer \SI{20}{\milli\second} beträgt.
  \item Bestimme \(Z\), Betrag und Winkel für \(R=\SI{8}{\ohm}\) und
        \(L=\SI{25}{\milli\henry}\) bei \(f=\SI{50}{\hertz}\).
  \item Ein symmetrischer Verbraucher nimmt bei \SI{400}{\volt} und
        \SI{25}{\ampere} einen Leistungsfaktor von \(0{,}82\) auf.
        Berechne \(P\), \(S\) und den Betrag von \(Q\).
  \item Erkläre, weshalb die Leiter-Erde-Spannung eines gesunden Leiters
        in einem isolierten Mittelspannungsnetz bei einem Erdschluss
        deutlich steigen kann, obwohl die Leiter-Leiter-Spannungen
        näherungsweise unverändert bleiben.
\end{aufgaben}

\section{Kurzlösungen}

\begin{enumerate}[label=\arabic*.]
  \item \(U=\hat U/\sqrt{2}\approx\SI{230}{\volt}\);
        \(f=1/T=\SI{50}{\hertz}\).
  \item \(\omega=2\pi f\approx\SI{314.16}{\radian\per\second}\),
        \(X_L=\omega L\approx\SI{7.85}{\ohm}\),
        \(Z=8+\jj7.85\ \si{\ohm}\),
        \(|Z|\approx\SI{11.21}{\ohm}\), Winkel etwa \(44.5^\circ\).
  \item \(S=\sqrt{3}UI\approx\SI{17.32}{\kilo\volt\ampere}\);
        \(P=S\cdot0.82\approx\SI{14.20}{\kilo\watt}\);
        \(|Q|=S\sqrt{1-0.82^2}\approx\SI{9.92}{\kilo var}\).
  \item Beim Erdschluss verschiebt sich der Sternpunkt gegen Erde. Die
        Leiter-Leiter-Spannungen werden durch die Verkettung der
        Leiterpotentiale bestimmt und können annähernd gleich bleiben,
        während die Leiter-Erde-Spannungen der gesunden Phasen steigen.
\end{enumerate}

\chapter{Netzgrößen, Leistung und Kurzschlussleistung}
\label{chap:netzgroessen}

\begin{lernziel}
Nach diesem Kapitel kannst du:
\begin{itemize}
  \item Netz-Nennspannung, Leiter-Leiter-Spannung und Strangspannung sicher unterscheiden,
  \item Scheinleistung und Nennstrom eines Transformators beziehungsweise Netzes berechnen,
  \item den Zusammenhang zwischen Kurzschlussleistung, Netzimpedanz und prospektivem Kurzschlussstrom herleiten,
  \item die Grenzen einer vereinfachten Kurzschlussrechnung benennen,
  \item Ergebnisse auf Plausibilität und Einheiten prüfen.
\end{itemize}
\end{lernziel}

\section{Nennwerte und tatsächliche Betriebsgrößen}

Die Nennspannung \(U_\mathrm{n}\) ist eine festgelegte Kenngröße zur Bezeichnung
eines Netzes oder Betriebsmittels. Sie ist nicht automatisch identisch mit der
Spannung, die zu jedem Zeitpunkt an jedem Knoten gemessen wird. Die tatsächliche
Betriebsspannung hängt unter anderem von Lastfluss, Netzregelung, Leitungsimpedanz
und Betriebszustand ab.

Bei Drehstromnetzen bezeichnet die übliche Netzspannungsangabe die Spannung
zwischen zwei Außenleitern. Für ein symmetrisches Netz gilt:
\[
U_\mathrm{LL}=\sqrt{3}\,U_\mathrm{Str}.
\]
Bei einem Netz mit \(U_\mathrm{LL}=\SI{20}{\kilo\volt}\) beträgt die
Strangspannung im symmetrischen Zustand daher
\[
U_\mathrm{Str}=\frac{\SI{20}{\kilo\volt}}{\sqrt{3}}
\approx \SI{11.55}{\kilo\volt}.
\]
Diese Beziehung darf nicht ohne Weiteres auf unsymmetrische Fehlerzustände
mit Sternpunktverschiebung übertragen werden.

\section{Drehstromleistung und Bemessungsstrom}

Für ein symmetrisches Drehstromsystem ist die Scheinleistung
\begin{equation}
  S=\sqrt{3}\,U_\mathrm{LL} I_\mathrm{L}.
  \label{eq:netz_scheinleistung}
\end{equation}
Daraus folgt für den Leiterstrom:
\begin{equation}
  I_\mathrm{L}=\frac{S}{\sqrt{3}\,U_\mathrm{LL}}.
  \label{eq:netz_nennstrom}
\end{equation}
Die Gleichung ist beispielsweise für die Plausibilisierung von
Transformator-Bemessungsströmen hilfreich. Sie sagt allein jedoch noch nichts
über zulässige Dauerbelastung, Kühlung, Umgebungstemperatur oder Schutzkonzept
aus.

\begin{rechenbeispiel}
Ein Transformator hat eine Bemessungsscheinleistung von
\(\SI{40}{\mega\volt\ampere}\) und eine Oberspannung von
\(\SI{110}{\kilo\volt}\). Der Bemessungsstrom auf der Oberspannungsseite ist
\[
I_\mathrm{n,OS}=
\frac{\SI{40}{\mega\volt\ampere}}
{\sqrt{3}\cdot\SI{110}{\kilo\volt}}
\approx \SI{210}{\ampere}.
\]
Auf einer Unterspannungsseite mit \(\SI{20}{\kilo\volt}\) ergibt sich bei
derselben Bemessungsscheinleistung:
\[
I_\mathrm{n,US}=
\frac{\SI{40}{\mega\volt\ampere}}
{\sqrt{3}\cdot\SI{20}{\kilo\volt}}
\approx \SI{1.155}{\kilo\ampere}.
\]
Die niedrigere Spannung führt bei gleicher Leistung zu einem höheren Strom.
Das ist einer der Gründe, weshalb Übertragungsnetze hohe Spannungen verwenden.
\end{rechenbeispiel}

\section{Netzimpedanz und Kurzschlussstrom}

Ein Kurzschluss führt zu einem hohen Strom, weil die Impedanz zwischen
speisender Quelle und Fehlerort stark kleiner ist als die Impedanz im
Normalbetrieb mit angeschlossener Last. Für eine erste einphasige Näherung
wird das Netz durch eine ideale Spannungsquelle hinter einer Ersatzimpedanz
\(Z_k\) modelliert.

Für einen Fehler an einer Effektivwertquelle gilt im einfachen Thevenin-Modell:
\begin{equation}
  \phasor I_k=\frac{\phasor U_\mathrm{Th}}{Z_k}.
  \label{eq:thevenin_ik}
\end{equation}
Der Betrag des Stroms ist damit
\[
I_k=\frac{U_\mathrm{Th}}{|Z_k|}.
\]
Wichtig: Spannung und Impedanz müssen auf dieselbe Ersatzschaltung und
Spannungsebene bezogen sein. In einer dreiphasigen Kurzschlussrechnung ist
bei Verwendung der Leiter-Leiter-Spannung die Strangspannung
\(U_\mathrm{LL}/\sqrt{3}\) die treibende Spannung des einphasigen
Ersatzmodells.

\section{Kurzschlussleistung}

Die Kurzschlussleistung ist eine Kenngröße für die elektrische Stärke eines
Netzknotens. In einer üblichen dreiphasigen Definition wird sie aus der
Nennspannung und dem anfänglichen symmetrischen Kurzschlusswechselstrom
gebildet:
\begin{equation}
  S_k''=\sqrt{3}\,U_\mathrm{n}\,I_k''.
  \label{eq:sk_definition}
\end{equation}
Daraus folgt:
\begin{equation}
  I_k''=\frac{S_k''}{\sqrt{3}\,U_\mathrm{n}}.
  \label{eq:ik_from_sk}
\end{equation}
Die doppelte Strichnotation \(I_k''\) wird in der Kurzschlussberechnung
häufig für den Anfangs-Kurzschlusswechselstrom verwendet. In konkreten
Normen und Rechenverfahren ist die jeweils definierte Stromgröße zu beachten.

\begin{merke}
Die Kurzschlussleistung \(S_k''\) ist keine während des Fehlers tatsächlich
in Wärme umgesetzte Wirkleistung. Sie ist eine aus Spannung und
Kurzschlussstrom gebildete Kenngröße des Netzknotens.
\end{merke}

\begin{rechenbeispiel}
An einem \(\SI{20}{\kilo\volt}\)-Netzknoten sei die dreiphasige
Kurzschlussleistung \(S_k''=\SI{250}{\mega\volt\ampere}\) angegeben.
Der anfängliche symmetrische Kurzschlussstrom beträgt näherungsweise
\[
I_k''=\frac{\SI{250}{\mega\volt\ampere}}
{\sqrt{3}\cdot\SI{20}{\kilo\volt}}
\approx \SI{7.22}{\kilo\ampere}.
\]
Die Größenordnung ist plausibel: Eine Kurzschlussleistung von
\(\SI{250}{\mega\volt\ampere}\) bei nur \(\SI{20}{\kilo\volt}\) führt zu
einem Strom im Bereich mehrerer Kiloampere.
\end{rechenbeispiel}

\section{Abschätzung der Netzimpedanz aus der Kurzschlussleistung}

Für eine erste Näherung eines symmetrischen dreipoligen Kurzschlusses kann
der Betrag der äquivalenten Netzimpedanz auf der betrachteten Spannungsebene
abgeschätzt werden durch
\begin{equation}
  |Z_k|\approx\frac{U_\mathrm{n}^2}{S_k''}.
  \label{eq:zk_from_sk}
\end{equation}
Dabei ist \(U_\mathrm{n}\) die Leiter-Leiter-Spannung. Die Formel ergibt sich
aus dem einphasigen Ersatzmodell mit Strangspannung \(U_\mathrm{n}/\sqrt{3}\)
und dem Zusammenhang \(S_k''=\sqrt{3}U_\mathrm{n}I_k''\):
\[
I_k''\approx\frac{U_\mathrm{n}/\sqrt{3}}{|Z_k|}
\quad\Longrightarrow\quad
S_k''\approx\frac{U_\mathrm{n}^2}{|Z_k|}.
\]
Die Beziehung liefert zunächst nur den Betrag. Für die Aufteilung in
\(R_k\) und \(X_k\) ist zusätzliche Information über das Verhältnis \(R/X\)
oder eine detailliertere Netzberechnung nötig.

\begin{rechenbeispiel}
Für den zuvor betrachteten Netzknoten gilt:
\[
|Z_k|\approx
\frac{(\SI{20}{\kilo\volt})^2}{\SI{250}{\mega\volt\ampere}}
=\SI{1.6}{\ohm}.
\]
Die Einheit lässt sich direkt prüfen:
\(\si{\volt\squared}/\si{\volt\ampere}=\si{\ohm}\).
Der Wert ist eine auf die \(\SI{20}{\kilo\volt}\)-Ebene bezogene
Ersatzimpedanz für die vereinfachte dreipolige Kurzschlussbetrachtung.
\end{rechenbeispiel}

\section{Was die einfache Rechnung nicht abbildet}

Die Formeln oben sind nützlich für erste Abschätzungen, ersetzen aber keine
vollständige Kurzschlussberechnung. Je nach Aufgabe sind unter anderem zu
berücksichtigen:
\begin{itemize}
  \item Impedanzen von Transformatoren, Leitungen und weiteren Quellen,
  \item Beiträge von Motoren, Generatoren und leistungselektronischen Anlagen,
  \item das Verhältnis \(R/X\) und der Gleichstromanteil des Kurzschlussstroms,
  \item der zeitliche Verlauf und die Unterscheidung zwischen Anfangs-,
        Ausschalt- und stationärem Kurzschlussstrom,
  \item Fehlerart, Fehlerimpedanz, Erdungsart und Nullsystemimpedanz,
  \item der verwendete Berechnungsstandard und dessen Spannungsfaktoren.
\end{itemize}

Ein besonders wichtiger Punkt: Die Kurzschlussleistung allein beschreibt
nicht alle Fehlerarten. Für einen Erdschluss oder einen zweipoligen Fehler
müssen die jeweiligen Sequenzimpedanzen und Netzerdungsverhältnisse
berücksichtigt werden.

\section{Plausibilitätsprüfungen}

Vor jeder Weiterverwendung sollte man mindestens prüfen:
\begin{enumerate}
  \item Stimmen die Spannungsebene und die Bezugsseite?
  \item Wurde Leiter-Leiter-Spannung oder Strangspannung eingesetzt?
  \item Ist der berechnete Strom in der erwarteten Größenordnung?
  \item Ist \(S_k''=\sqrt{3}U_nI_k''\) mit den eingesetzten Werten erfüllt?
  \item Ist klar, ob ein Anfangs-, Ausschalt- oder Dauerstrom gemeint ist?
  \item Handelt es sich um einen dreipoligen Kurzschluss oder einen Erdschluss?
\end{enumerate}

\section{Übungsaufgaben}

\begin{aufgaben}
  \item Ein \(\SI{63}{\mega\volt\ampere}\)-Transformator hat auf der
        \(\SI{110}{\kilo\volt}\)-Seite seine Bemessungsspannung.
        Berechne den Bemessungsstrom auf dieser Seite.
  \item Ein \(\SI{10}{\kilo\volt}\)-Netzknoten hat eine angegebene
        dreiphasige Kurzschlussleistung von \(\SI{180}{\mega\volt\ampere}\).
        Berechne den anfänglichen symmetrischen Kurzschlussstrom.
  \item Bestimme aus Aufgabe 2 den Betrag der vereinfachten
        Thevenin-Ersatzimpedanz auf der \(\SI{10}{\kilo\volt}\)-Ebene.
  \item Ein Techniker setzt in einem dreiphasigen Kurzschlussmodell
        \(U_\mathrm{n}\) direkt als Strangspannung ein. Erkläre den Fehler
        und die Auswirkung auf das Ergebnis.
  \item Warum reicht die Kurzschlussleistung eines Netzknotens allein
        nicht aus, um einen einpoligen Erdschlussstrom zu bestimmen?
\end{aufgaben}

\section{Kurzlösungen}

\begin{enumerate}[label=\arabic*.]
  \item \(I_n=63/(\sqrt{3}\cdot110)\,\si{\mega\volt\ampere\per\kilo\volt}
        \approx\SI{331}{\ampere}\).
  \item \(I_k''=180/(\sqrt{3}\cdot10)\,\si{\mega\volt\ampere\per\kilo\volt}
        \approx\SI{10.39}{\kilo\ampere}\).
  \item \(|Z_k|\approx U_n^2/S_k''=(10\,\si{\kilo\volt})^2/
        180\,\si{\mega\volt\ampere}\approx\SI{0.556}{\ohm}\).
  \item Die treibende Spannung des einphasigen Ersatzmodells ist
        \(U_n/\sqrt{3}\). Eine Verwechslung um den Faktor \(\sqrt{3}\)
        führt entsprechend zu einem falschen Strom.
  \item Dafür werden insbesondere die Nullsystemimpedanz und die
        Sternpunktbehandlung benötigt; diese werden in späteren Kapiteln
        eingeführt.
\end{enumerate}
\chapter{Per-Unit-System und Impedanzumrechnung}
\label{chap:perunit}

\begin{lernziel}
Nach diesem Kapitel können Sie Bezugsgrößen eines Drehstromnetzes festlegen,
Impedanzen auf eine andere Spannungsebene beziehen, aus der
Kurzschlussspannung eines Transformators dessen Impedanz bestimmen und
mehrere Netzkomponenten auf eine gemeinsame Per-Unit-Basis umrechnen. Sie
können außerdem die Grenzen einer überschlägigen Rechnung benennen.
\end{lernziel}

\section{Warum Bezugsgrößen helfen}

Ein Netzmodell kann eine Einspeisung, mehrere Transformatoren und Leitungen
verschiedener Spannungsebenen enthalten. In einem rein ohmsch-induktiven
Ersatzschaltbild müssen alle Impedanzen, die addiert werden, auf dieselbe
Spannungsebene bezogen sein. Das klassische Umrechnen über jedes
Transformatorübersetzungsverhältnis ist korrekt, wird bei mehreren Ebenen aber
unübersichtlich. Das Per-Unit-System (p.u.) beschreibt eine Größe als Anteil
einer festgelegten Bezugsgröße.

Für ein symmetrisches Drehstromsystem legen wir üblicherweise eine
Bezugs-Scheinleistung $S_B$ und eine Leiter-Leiter-Bezugsspannung $U_B$ fest.
Daraus ergeben sich die Bezugsgrößen
\begin{align}
  I_B &= \frac{S_B}{\sqrt{3}\,U_B},\\
  Z_B &= \frac{U_B^2}{S_B},\\
  Y_B &= \frac{1}{Z_B}=\frac{S_B}{U_B^2}.
\end{align}
Dabei sind $S_B$ die gesamte dreiphasige Scheinleistung und $U_B$ die
Leiter-Leiter-Spannung. Mit $U_B$ in Volt und $S_B$ in Voltampere erhält man
$Z_B$ in Ohm. In der Praxis dürfen konsistente Präfixe wie kV und MVA
verwendet werden, sofern das Ergebnis dimensionsrichtig gebildet wird.

Die normierte Impedanz lautet
\begin{equation}
  z_{\mathrm{pu}}=\frac{Z}{Z_B}=Z\frac{S_B}{U_B^2}.
  \label{eq:zpu}
\end{equation}
Analog gilt $i_{\mathrm{pu}}=I/I_B$ und $u_{\mathrm{pu}}=U/U_B$.
Bei Wechselstromrechnung sind Impedanzen und Spannungen komplexe Größen; die
Bezugsgrößen selbst sind reelle positive Beträge.

\begin{merke}
Die Formel $Z_B=U_B^2/S_B$ gilt in dieser Form für ein konsistentes
Drehstromsystem mit Leiter-Leiter-Bezugsspannung und dreiphasiger
Bezugsleistung. Vermischen Sie diese Größen nicht mit einer einphasigen
Rechnung, ohne die Definitionen entsprechend anzupassen.
\end{merke}

\section{Impedanz auf eine andere Spannungsebene beziehen}

Für einen idealen Transformator mit Übersetzungsverhältnis
\begin{equation}
  a=\frac{U_1}{U_2}
\end{equation}
folgt aus der Erhaltung der Leistung und dem Übersetzungsverhältnis der
Spannungen die auf Seite 1 bezogene Impedanz einer Komponente auf Seite 2:
\begin{equation}
  Z_2'=a^2 Z_2.
  \label{eq:impedanz_beziehen}
\end{equation}
In Gegenrichtung gilt $Z_1'=Z_1/a^2$. Die Regel gilt für die gesamte komplexe
Impedanz, also sowohl für den Widerstands- als auch für den Reaktanzanteil:
\begin{equation}
  R'=a^2R,\qquad X'=a^2X.
\end{equation}

\begin{rechenbeispiel}
Eine Niederspannungs-Komponente habe $Z_2=(0{,}08+j0{,}06)\,\Omega$. Ein
idealer Transformator setze $20\,\mathrm{kV}$ auf $0{,}4\,\mathrm{kV}$ um.
Dann ist $a=50$ und damit
\[
 Z_2'=50^2(0{,}08+j0{,}06)\,\Omega
     =(200+j150)\,\Omega.
\]
Der Betrag der Impedanz wird mit $a^2=2500$ multipliziert. Das ist kein
physikalisches Anwachsen der Komponente, sondern die Darstellung derselben
Wirkung auf der Hochspannungsseite des Ersatzschaltbilds.
\end{rechenbeispiel}

\section{Per-Unit-Umrechnung zwischen Bezugsbasen}

Eine physikalische Impedanz bleibt gleich, ihr p.u.-Wert hängt aber von der
gewählten Bezugsleistung und Bezugsspannung ab. Ist eine Impedanz auf der
alten Basis $(S_{B,\mathrm{alt}},U_{B,\mathrm{alt}})$ als
$z_{\mathrm{pu,alt}}$ bekannt, ergibt sich auf einer neuen Basis
\begin{equation}
 z_{\mathrm{pu,neu}}=z_{\mathrm{pu,alt}}
 \frac{S_{B,\mathrm{neu}}}{S_{B,\mathrm{alt}}}
 \left(\frac{U_{B,\mathrm{alt}}}{U_{B,\mathrm{neu}}}\right)^2.
 \label{eq:pu_basiswechsel}
\end{equation}
Die Formel folgt direkt daraus, dass $Z=z_{\mathrm{pu}}U_B^2/S_B$ auf beiden
Basen dieselbe reale Impedanz beschreiben muss.

Wählt man die Bezugsspannungen auf beiden Seiten eines Transformators im
Verhältnis seiner Nennspannungen, vereinfacht sich die Umrechnung über den
Transformator hinweg: Die p.u.-Impedanz bleibt bei gleicher Bezugsleistung
auf beiden Seiten gleich. Genau das macht die Methode für mehrstufige Netze
so nützlich.

\begin{merke}
Beim Basiswechsel darf der Spannungsfaktor in
\cref{eq:pu_basiswechsel} nicht weggelassen werden, wenn die alte und neue
Bezugsspannung nicht identisch sind. Die häufig verwendete Kurzform
$z_{\mathrm{pu,neu}}=z_{\mathrm{pu,alt}}S_{B,\mathrm{neu}}/S_{B,\mathrm{alt}}$
setzt voraus, dass die Bezugsspannung unverändert bleibt oder dass die
Bezugsbasen entlang eines Transformators passend gewählt sind.
\end{merke}

\section{Transformatorimpedanz aus Typenschilddaten}

Auf dem Transformator-Typenschild findet sich häufig die relative
Kurzschlussspannung $u_k$ in Prozent. Sie ist die Spannung, die bei
kurzgeschlossener Gegenseite erforderlich ist, damit bei Nennfrequenz der
Nennstrom fließt, angegeben als Anteil der Nennspannung. Für das übliche
vereinfachte Ersatzschaltbild ergibt sich der Betrag der Impedanz bezogen auf
die Nennbasis des Transformators zu
\begin{equation}
  z_{T,\mathrm{pu}}\approx\frac{u_k[\%]}{100}.
  \label{eq:uk_pu}
\end{equation}
Die Näherung bezieht sich auf den Impedanzbetrag. $u_k$ allein liefert nicht
die Aufteilung in Widerstand und Reaktanz. Wenn die Kupferverluste $P_k$
(Nennstrom-Kurzschlussverluste) bekannt sind, lässt sich diese Aufteilung
näherungsweise ergänzen.

Für die Bezugsleistung $S_N$ und die Leiter-Leiter-Nennspannung $U_N$ ist der
Nennstrom
\[
 I_N=\frac{S_N}{\sqrt{3}U_N}.
\]
Der Impedanzbetrag auf der jeweiligen Seite beträgt
\begin{equation}
 |Z_T|=\frac{u_k}{100}\frac{U_N^2}{S_N},
 \label{eq:zt_ohm}
\end{equation}
sofern $u_k$ als dimensionsloser Wert eingesetzt wird. Für die Umrechnung auf
eine andere Bezugsbasis wird anschließend
\cref{eq:pu_basiswechsel} verwendet.

\begin{rechenbeispiel}
Ein Transformator hat $S_N=\SI{10}{\mega\volt\ampere}$,
$U_N=\SI{20}{\kilo\volt}$ auf der betrachteten Seite und $u_k=\SI{8}{\percent}$.
Dann gilt
\[
 Z_B=\frac{(20\,\mathrm{kV})^2}{10\,\mathrm{MVA}}=40\,\Omega,
 \qquad |Z_T|\approx0{,}08\cdot40\,\Omega=3{,}2\,\Omega.
\]
Das ist der Betrag der auf diese Seite bezogenen Kurzschlussimpedanz, nicht
automatisch ein reiner Blindwiderstand von $j3{,}2\,\Omega$.
\end{rechenbeispiel}

\section{Netzimpedanz aus Kurzschlussleistung}

Ist am Einspeisepunkt die dreiphasige Anfangs-Kurzschlussleistung
$S_k''$ bekannt, kann der Betrag der äquivalenten Netzimpedanz überschlägig
bestimmt werden:
\begin{equation}
 |Z_Q|\approx\frac{U_n^2}{S_k''}.
 \label{eq:zquelle}
\end{equation}
Auf einer p.u.-Basis mit $U_B=U_n$ folgt unmittelbar
\begin{equation}
 |z_{Q,\mathrm{pu}}|\approx\frac{S_B}{S_k''}.
 \label{eq:zquellepu}
\end{equation}
Die Kurzschlussleistung allein bestimmt wiederum nicht den Winkel der
Impedanz. Für eine belastbare Kurzschlussrechnung werden, soweit erforderlich,
auch das Verhältnis $R/X$, die Netzbetriebsbedingungen, die Beiträge von
Motoren oder dezentralen Einspeisern und die anzuwendenden Rechenregeln
berücksichtigt.

\section{Durchgängiges Beispiel: Quelle, Transformator und Kabel}

Gegeben sei ein vereinfachtes Netz mit folgenden Daten:
\begin{itemize}
 \item Einspeisung auf der Mittelspannungsseite: $U_n=\SI{20}{\kilo\volt}$,
       $S_k''=\SI{250}{\mega\volt\ampere}$,
 \item Transformator: $\SI{20}{\kilo\volt}/\SI{0.4}{\kilo\volt}$,
       $S_N=\SI{1}{\mega\volt\ampere}$, $u_k=\SI{6}{\percent}$,
 \item Niederspannungskabel: $Z_K=(0{,}015+j0{,}010)\,\Omega$,
       auf die betrachtete Strecke bezogen.
\end{itemize}

Wir wählen $S_B=\SI{1}{\mega\volt\ampere}$ als gemeinsame
Bezugsleistung. Die Bezugsspannungen sind $\SI{20}{\kilo\volt}$ auf der
Mittelspannungsseite und $\SI{0.4}{\kilo\volt}$ auf der
Niederspannungsseite. Sie stehen im selben Verhältnis wie die
Transformator-Nennspannungen.

\paragraph{1. Einspeisung.}
Mit \cref{eq:zquellepu} folgt
\[
 |z_{Q,\mathrm{pu}}|\approx\frac{1}{250}=0{,}004.
\]
Für die Beispielrechnung nehmen wir vereinfachend eine rein induktive
Quellenimpedanz $z_Q\approx j0{,}004$ an. Das ist eine Modellannahme, keine
Information, die sich allein aus $S_k''$ ableiten lässt.

\paragraph{2. Transformator.}
Da die Bezugsleistung gleich der Transformator-Nennleistung ist und die
Bezugsspannungen den Nennspannungen folgen, gilt
\[
 z_{T,\mathrm{pu}}\approx j0{,}06,
\]
unter der zusätzlichen vereinfachenden Annahme, dass die Transformatorimpedanz
vollständig reaktiv ist. Aus $u_k$ allein wäre diese Annahme nicht beweisbar.

\paragraph{3. Kabel.}
Auf der Niederspannungsseite ist
\[
 Z_{B,\mathrm{NS}}=\frac{(0{,}4\,\mathrm{kV})^2}{1\,\mathrm{MVA}}
 =0{,}16\,\Omega.
\]
Somit
\[
 z_{K,\mathrm{pu}}=\frac{0{,}015+j0{,}010}{0{,}16}
 =0{,}09375+j0{,}0625.
\]

\paragraph{4. Gesamtimpedanz und Fehlerstrom.}
Alle Impedanzen liegen nun auf derselben p.u.-Basis und können addiert werden:
\begin{align*}
 z_{\Sigma,\mathrm{pu}}&=z_Q+z_T+z_K\\
 &\approx0{,}09375+j(0{,}004+0{,}06+0{,}0625)\\
 &=0{,}09375+j0{,}1265.
\end{align*}
Der Betrag ist $|z_{\Sigma,\mathrm{pu}}|\approx0{,}1574$. Der Bezugsstrom auf
der Niederspannungsseite beträgt
\[
 I_{B,\mathrm{NS}}=\frac{1\,\mathrm{MVA}}{\sqrt{3}\cdot0{,}4\,\mathrm{kV}}
 \approx\SI{1.443}{\kilo\ampere}.
\]
Für einen dreipoligen Fehler am Kabelende ergibt die vereinfachte Rechnung
\[
 I_k\approx\frac{I_B}{|z_{\Sigma,\mathrm{pu}}|}
 \approx\SI{9.17}{\kilo\ampere}.
\]

\begin{sicherheitshinweis}
Dieses Beispiel illustriert die Rechenmethode, ist aber keine vollständige
Kurzschlussberechnung nach einer anzuwendenden Norm. Es vernachlässigt unter
anderem den ohmschen Anteil der Quellen- und Transformatorimpedanz,
Spannungsfaktoren, Toleranzen, Betriebszustände und mögliche weitere
Kurzschlussstrombeiträge. Für die Auslegung von Schaltgeräten, Schutzorganen
und Betriebsmitteln müssen die maßgeblichen Randbedingungen und aktuellen
Regelwerke berücksichtigt werden.
\end{sicherheitshinweis}

\section{Typische Fehler bei der Umrechnung}
\begin{itemize}
 \item Leiter-Leiter-Spannung und Phasenspannung werden in einer Rechnung
       vermischt.
 \item Eine Impedanz wird ohne Quadrat des Übersetzungsverhältnisses auf eine
       andere Spannungsebene bezogen.
 \item Bei einem Basiswechsel wird der Faktor für die Bezugsleistung oder
       Bezugsspannung vergessen.
 \item Der Prozentwert $u_k=6\,\%$ wird als $6$ statt als $0{,}06$ eingesetzt.
 \item Aus $u_k$ oder $S_k''$ wird fälschlich ein vollständiger komplexer
       Impedanzwert einschließlich Winkel abgeleitet.
 \item Impedanzen in Ohm von unterschiedlichen Spannungsebenen werden direkt
       addiert, ohne sie vorher umzurechnen.
\end{itemize}

\begin{aufgaben}
 \item Bestimmen Sie $Z_B$ für $S_B=\SI{5}{\mega\volt\ampere}$ und
       $U_B=\SI{10}{\kilo\volt}$.
 \item Ein Transformator hat $S_N=\SI{5}{\mega\volt\ampere}$,
       $U_N=\SI{10}{\kilo\volt}$ und $u_k=\SI{7}{\percent}$. Bestimmen Sie
       den Impedanzbetrag auf dieser Seite in Ohm.
 \item Eine Impedanz beträgt auf einer Basis $z_{\mathrm{pu,alt}}=0{,}08$.
       Die Bezugsspannung bleibt gleich; die Bezugsleistung steigt von
       $\SI{10}{\mega\volt\ampere}$ auf $\SI{20}{\mega\volt\ampere}$.
       Wie groß ist $z_{\mathrm{pu,neu}}$?
 \item Warum reicht die Angabe $u_k=\SI{6}{\percent}$ nicht aus, um den
       vollständigen komplexen Transformatorimpedanzwert zu bestimmen?
 \item Welche Bedingung muss erfüllt sein, bevor Impedanzen in Ohm direkt
       addiert werden dürfen?
\end{aufgaben}

\section*{Kurzlösungen}
\begin{enumerate}[label=\arabic*.]
 \item $Z_B=(10\,\mathrm{kV})^2/(5\,\mathrm{MVA})=20\,\Omega$.
 \item $Z_B=20\,\Omega$ und $|Z_T|\approx0{,}07\cdot20=1{,}4\,\Omega$.
 \item $z_{\mathrm{pu,neu}}=0{,}08\cdot20/10=0{,}16$.
 \item $u_k$ liefert den Impedanzbetrag; für den Winkel bzw. die Aufteilung in
       $R$ und $X$ braucht man weitere Angaben, etwa Kurzschlussverluste.
 \item Sie müssen auf dieselbe Spannungsebene bezogen sein und sich auf
       kompatible Ersatzschaltbild- und Frequenzannahmen beziehen.
\end{enumerate}


\part{Betriebsmittel und Ersatzschaltbilder}
\chapter*{Transformatoren und Schaltgruppen}
\label{chap:04-transformatoren}

Dieses Kapitel ist als Bestandteil des modularen Lehrskripts vorgesehen.
Inhalte werden ergänzt, nachdem die fachlichen Grundlagen und die
Quellenlage für dieses Thema systematisch ausgearbeitet wurden.

\begin{center}
\begin{minipage}{0.8\textwidth}
\noindent
\begin{tikzpicture}[transform shape]
	% Paths, nodes and wires:
	\draw (2, 4) to[european inductor, l={$L_{s1}=\sigma_1 L_1$}] (4, 4);
	\draw (0, 4) to[generic, l={$R_1$}] (2, 4);
	\draw (5, 2) to[european inductor, l_={$L_{h1}$}, label distance=0.03cm] (5, 4);
	\draw (10, 4) to[generic, l={$R_2$}] (12, 4);
	\draw (8, 4) to[european inductor, l={$L_{s2}=\sigma_2 L_2$}] (10, 4);
	\draw (4, 2) to[generic, l={$R_{FE}$}] (4, 4);
	\node[transformer core, cute, xscale=0.95, yscale=0.95] at (7.003, 3.003){};
	\draw (6.005, 4) |- (5, 4);
	\draw (6, 2) |- (5, 2) |- (4, 2) -- (0, 2);
	\draw (4, 4) -- (5, 4);
	\draw (8, 2.005) -| (12, 2);
	\draw[line width=1pt, -latex] (0, 3.75) -- (0, 2.25)   node[midway, right=2pt] {$\underline{U}_1$};
	\draw[line width=1pt, -latex] (12, 3.75) -- (12, 2.25) node[midway, right=2pt] {$\underline{U}_2$};
	\draw[line width=1pt, -latex] (6, 3.75) -- (6, 2.25)   ;%node[midway, right=2pt] {$\underline{U}_{1h}$};
	\draw[line width=1pt, -latex] (8, 3.75) -- (8, 2.25)   node[midway, right=2pt] {$\underline{U}_{2h} = \frac{N_2}{N_1} \, \underline{U}_{1h} $};
    \node[above=3pt] at (5.8, 4) {$L_{h1} = (1-\sigma_1) L_1 $};
\end{tikzpicture}
\begin{tikzpicture}[transform shape]
	% Paths, nodes and wires:
	\draw (2, 4) to[european inductor, l={$L_{s1}=\sigma_1 L_1$}] (4, 4);
	\draw (0, 4) to[generic, l={$R_1$}] (2, 4);
	\draw (5, 2) to[european inductor, l_={$L_{h1}$}, label distance=0.03cm] (5, 4);
	\draw (10, 4) to[generic, l={$R_2$}] (12, 4);
	\draw (8, 4) to[european inductor, l={$L_{s2}=\sigma_2 L_2$}] (10, 4);
	\draw (4, 2) to[generic, l={$R_{FE}$}] (4, 4);
	\node[transformer core, cute, xscale=0.95, yscale=0.95] at (7.003, 3.003){};
	\draw (6.005, 4) |- (5, 4);
	\draw (6, 2) |- (5, 2) |- (4, 2) -- (0, 2);
	\draw (4, 4) -- (5, 4);
	\draw (8, 2.005) -| (12, 2);
	\draw[line width=1pt, -latex] (0, 3.75) -- (0, 2.25)   node[midway, right=2pt] {$\underline{U}_1$};
	\draw[line width=1pt, -latex] (12, 3.75) -- (12, 2.25) node[midway, right=2pt] {$\underline{U}_2$};
	\draw[line width=1pt, -latex] (6, 3.75) -- (6, 2.25)   ;%node[midway, right=2pt] {$\underline{U}_{1h}$};
	\draw[line width=1pt, -latex] (8, 3.75) -- (8, 2.25)   node[midway, right=2pt] {$\underline{U}_{2h} = \frac{N_2}{N_1} \, \underline{U}_{1h} $};
    \node[above=3pt] at (5.8, 4) {$L_{h1} = (1-\sigma_1) L_1 $};
\end{tikzpicture}
\par\bigskip
\noindent
\begin{tikzpicture}[transform shape]
	% Paths, nodes and wires:
	\draw (2, 4) to[european inductor, l={$L_{s1}=\sigma_1 L_1$}] (4, 4);
	\draw (0, 4) to[generic, l={$R_1$}] (2, 4);
	\draw (5, 2) to[european inductor, l_={$L_{h1}$}, label distance=0.03cm] (5, 4);
	\draw (8.005, 4) to[generic, l={$R_2'$}, i^<=$i_2'$] (10.005, 4);
	\draw (6.005, 4) to[european inductor, l={$L_{s2}'=\sigma_2 L_2'$}] (8.005, 4);
	\draw (4, 2) to[generic, l={$R_{FE}$}] (4, 4);
	\draw (6.005, 4) |- (5, 4);
	\draw (10, 2) -- (5, 2) |- (4, 2) -- (0, 2);
	\draw (4, 4) -- (5, 4);
    \draw (6, 2) |- (5, 2) |- (4, 2) -- (0, 2);
	\draw (8, 2.005) -| (10, 2);
	\draw (2, 4) to[european inductor, l={$L_{s1}=\sigma_1 L_1$}] (4, 4);
	\draw (0, 4) to[generic, l={$R_1$}] (2, 4);
	\draw (5, 2) to[european inductor, l_={$L_{h1}$}, label distance=0.03cm] (5, 4);
	\draw (4, 2) to[generic, l={$R_{FE}$}] (4, 4);
	\node[transformer core, cute, xscale=0.95, yscale=0.95] at (10.997, 3.003){};
	\draw[line width=1pt, -latex] (0, 3.75) -- (0, 2.25)   node[midway, right=2pt] {$\underline{U}_1$};
	\draw[line width=1pt, -latex] (12, 3.75) -- (12, 2.25) node[midway, right=2pt] {$\underline{U}_2$};
	\draw[line width=1pt, -latex] (10, 3.75) -- (10, 2.25) node[midway, left=0pt] {$\frac{N_1}{N_2} \, \underline{U}_2$};
	\draw[line width=1pt, -latex] (6, 3.75) -- (6, 2.25)  node[midway, right=2pt]  {$\underline{U}_{1h} $};
\end{tikzpicture}
\par\bigskip
\noindent
\begin{tikzpicture}[transform shape]
	% Paths, nodes and wires:
	\draw (2, 4) to[european inductor, l={$L_{s1}=\sigma_1 L_1$}] (4, 4);
	\draw (0, 4) to[generic, l={$R_1$}] (2, 4);
	\draw (5, 2) to[european inductor, l_={$L_{h1}$}, label distance=0.03cm] (5, 4);
	\draw (8, 4) to[generic, l={$R_2'$}, i^<=$i_2'$]  (10, 4);
	\draw (6, 4) to[european inductor, l={$L_{s2}'$}] (8, 4);
	\draw (4, 2) to[generic, l={$R_{FE}$}] (4, 4);
	\draw (6, 4) |- (5, 4);
	\draw (10, 2) -- (5, 2) |- (4, 2) -- (0, 2);
	\draw (4, 4) -- (5, 4);
	\draw[line width=1pt, -latex] (0, 3.75) -- (0, 2.25) node[midway, right=2pt] {$\underline{U}_1$};
	\draw[line width=1pt, -latex] (10, 3.75) -- (10, 2.25) node[midway, right=2pt] {$\underline{U}_2'$};
	\draw[line width=1pt, -latex] (6, 3.75) -- (6, 2.25) node[midway, right=2pt] {$\underline{U}_{1h}$};
\end{tikzpicture}

\par\bigskip
\noindent
\begin{tikzpicture}[transform shape]
	% Paths, nodes and wires:
	\draw (4.5, 5)    to[european inductor, l={$X_{1\sigma}+X_{2\sigma}'$}] (6.5, 5);
	\draw (2.5, 5)    to[generic, l={$R_1+R_2'$}] (4.5, 5);
	\draw (1.5, 2.5)  to[european inductor, l={$2X_{1h}$}, label distance=0.03cm] (1.5, 4.5);
	\draw (2.5, 2.5)  to[generic, l_={$2R_{FE}$}] (2.5, 4.5);
	\draw (6.5, 2.5)  to[european inductor, l={$2X_{1h}$}, label distance=0.03cm] (6.5, 4.5);
	\draw (7.5, 2.5)  to[generic, l_={$2R_{FE}$}] (7.5, 4.5);
    \draw[line width=1pt, -latex] (0, 4.5) -- (0, 2.5) node[midway, left=2pt] {$\underline{U}_{1}$};
	\draw[line width=1pt, -latex] (9, 4.5) -- (9, 2.5) node[midway, right=2pt] {$\underline{U}_{2}'$};
	\draw (1.5, 2.5) -- (2.5, 2.5);
	\draw (1.5, 4.5) -- (2.5, 4.5);
	\draw (0, 5) -- (2.5, 5);
    \draw (0, 5) to[short, i^>={$I_1$}] (1,5);
	\draw (2, 4.5) -| (2, 5);
	\draw (6.5, 2.5) -- (7.5, 2.5);
	\draw (6.5, 4.5) -- (7.5, 4.5);
	\draw (7, 4.5) -| (7, 5);
	\draw (0, 2) -- (9, 2);
	\draw (2, 2.5) -| (2, 2);
	\draw (7, 2.5) -| (7, 2);
	\draw (6.5, 5) -- (9, 5);
    \draw (8, 5) to[short, i^<={$I_2'$}] (9,5);
\end{tikzpicture}


\begin{tikzpicture}[transform shape]
	\draw (1, 0.75) to[tmultiwire] (1, 0.25);
	\draw (1, 2.5) to[ioosource] (1, -0);
	\draw (1, 2.25) to[tmultiwire] (1, 1.75);
\end{tikzpicture}

\end{minipage}
\end{center}
$$ L_1 = L_{h1} + L_{s1} $$
$$ L_{h1} = (1-\sigma_1) L_1 $$

$$ M = \frac{N_1}{N_2} = k \, \sqrt{L_1 L_2} $$


\chapter{Transformatoren und ihre Ersatzschaltbilder}
\label{chap:transformatoren}

\section{Lernziele und Einordnung}

Transformatoren ermöglichen die verlustarme Übertragung elektrischer Energie
zwischen unterschiedlichen Spannungsebenen. Für die Netzberechnung sind sie
nicht nur als ideale Spannungsübersetzer zu betrachten. Wicklungswiderstände,
Streuflüsse, Magnetisierungsstrom, Eisenverluste und die Schaltgruppe bestimmen
das Betriebsverhalten und beeinflussen insbesondere Kurzschluss- und Erdschlussströme.

Nach diesem Kapitel sollen folgende Aufgaben gelöst werden können:
\begin{itemize}
  \item Übersetzungsverhältnis und Nennströme aus Typenschilddaten bestimmen;
  \item ein einphasiges Ersatzschaltbild eines Drehstromtransformators aufstellen;
  \item Leerlauf- und Kurzschlussversuch auswerten;
  \item aus der Kurzschlussspannung eine Ersatzimpedanz bestimmen;
  \item Wirk- und Blindanteil der Transformatorimpedanz unterscheiden;
  \item die Auswirkungen von Schaltgruppe und Neutralpunkterdung auf das
        Nullsystem qualitativ beurteilen.
\end{itemize}

\section{Idealer Transformator}

\subsection{Induktionsgesetz und Spannungsübersetzung}

Die elektromagnetische Induktion wird durch das Faradaysche
Induktionsgesetz beschrieben. In differentieller Form der
Maxwellgleichungen enthalten als
\[
    \nabla \times \vec{E}
    = -\frac{\partial \vec{B}}{\partial t}.
\]
Sie beschreibt, wie eine zeitlich veränderliche magnetische
Flussdichte $\vec{B}$ ein elektrisches Wirbelfeld erzeugt.
Durch Integration über eine von
einer Wicklung umschlossene Fläche folgt das Induktionsgesetz für
eine einzelne Windung. Für eine Wicklung mit $N$ Windungen ergibt
sich unter Berücksichtigung der gewählten Bezugsrichtungen
\[
    u(t) = N \frac{\mathrm d\Phi(t)}{\mathrm dt},
\]
wobei $\Phi$ der magnetische Fluss durch eine Windungsfläche ist.
Das Vorzeichen hängt von der Festlegung der Wicklungsrichtung bezüglich
der positiven Flussrichtung ab.

\subsection{Entstehung des magnetischen Flusses}

Der magnetische Fluss im Transformatorkern wird durch den Strom
in den Wicklungen hervorgerufen. Das Durchflutungsgesetz, oder
in vollständiger, differentieller Form, das
Ampère-Maxwell-Gesetz

\[
        \nabla \times \vec{H}
        = \vec{J} + \frac{\partial}{\partial t} \vec{D}
\]

beschreibt mit $\vec{H}$ die entstehende magnetische Feldstärke
aufgrund von elektrischer Stromdichte $\vec{J}$ und in seiner
Wirkung als Verschiebungsstrom bezeichneten Ableitung der elektrischen Flussdichte $\vec{D}$. Das
Gesetz besagt, dass sowohl elektrische Ströme als auch zeitlich veränderliche
elektrische Felder ein magnetisches Wirbelfeld hervorrufen.

Der Wicklungsstrom $N I$ bildet für einen Trafo die Durchflutung,
die nach diesem Gesetz mit einer mittleren
magnetischen Weglänge $l_\mathrm{m}$ dem entstehenden magnetischen Feld gleicht.

\[
    H l_\mathrm{m} = N I.
\]

Die magnetische Flussdichte ergibt sich aus der Materialgleichung
\[
    \vec{B} = \mu_0 \mu_\mathrm{r}\vec{H},
\]
wobei $\mu_0$ die magnetische Feldkonstante und $\mu_\mathrm{r}$
die relative Permeabilität des Kernmaterials bezeichnet. Der
magnetische Fluss durch den Kernquerschnitt $A_\mathrm{K}$ ist
\[
    \Phi = \int_{A_\mathrm{K}} \vec{B}\cdot \mathrm d\vec{A}.
\]
Bei näherungsweise homogener Flussdichte und einem konstanten
Kernquerschnitt vereinfacht sich dies zu
\[
    \Phi \approx B A_\mathrm{K}.
\]

Für den nachfolgend skizzierten Transformator ist die als $\underline{U}_1$
eingespeiste Spannung die Ursache für einen Strom $i_1$, der zu einem Fluss
$\Phi$ führt und eine der Ursache entgegenwirkende, induszierte Spannung
aufbaut, die das Treiben von Strom zu dem sich einstellenden Arbeitspunkt begrenzt.

\vspace{-1.1em}
\begin{center}
\begin{tikzpicture}[transform shape]
	% Paths, nodes and wires:
	\node[shape=rectangle, fill={rgb,255:red,120;green,150;blue,165}, draw, line width=1pt, minimum width=4.965cm, minimum height=5.715cm] at (3.49, 2.86){};
	\draw[line width=1pt, dash pattern={on 4pt off 2pt on 1pt off 2pt}, -latex] (3.5, 5) -- (5, 5) -- (5, 0.75) -- (2, 0.75) -- (2, 5) -- (3.75, 5) node[midway, above=2pt] {$\underline{\Phi}(t)$};
	\node[shape=rectangle, fill={rgb,255:red,255;green,255;blue,255}, draw, line width=1pt, minimum width=2.477cm, minimum height=3.715cm] at (3.496, 2.86){};
	\node[shape=rectangle, fill={rgb,255:red,209;green,209;blue,209}, draw, line width=1pt, minimum width=2.251cm, minimum height=3.215cm] at (1.597, 2.86){};
	\node[shape=circle, draw, line width=1pt, minimum width=0.465cm] at (2.404, 3.604){};
	\draw[line width=1pt] (2.227, 3.78) -- (2.581, 3.427);
	\draw[line width=1pt] (2.227, 3.427) -- (2.581, 3.78);
	\node[shape=circle, draw, line width=1pt, minimum width=0.465cm] at (0.802, 2.485){};
	\node[circ]  at (0.802, 2.485){};
	\node[shape=circle, draw, line width=1pt, minimum width=0.465cm] at (0.802, 3.985){};
	\node[circ]  at (0.802, 3.985){};
	\node[shape=circle, draw, line width=1pt, minimum width=0.465cm] at (0.802, 1.735){};
	\node[circ]  at (0.802, 1.735){};
	\node[shape=circle, draw, line width=1pt, minimum width=0.465cm] at (2.404, 2.854){};
	\draw[line width=1pt] (2.227, 3.03) -- (2.581, 2.677);
	\draw[line width=1pt] (2.227, 2.677) -- (2.581, 3.03);
	\node[shape=circle, draw, line width=1pt, minimum width=0.465cm] at (2.404, 2.104){};
	\draw[line width=1pt] (2.227, 2.28) -- (2.581, 1.927);
	\draw[line width=1pt] (2.227, 1.927) -- (2.581, 2.28);
	\node[shape=circle, draw, line width=1pt, minimum width=0.465cm] at (0.802, 3.235){};
	\node[circ]  at (0.802, 3.235){};
	\node[shape=rectangle, xscale=-1, yscale=-1, fill={rgb,255:red,209;green,209;blue,209}, draw, line width=1pt, minimum width=2.251cm, minimum height=3.215cm] at (5.37, 2.875){};
	\node[shape=circle, xscale=-1, yscale=-1, draw, line width=1pt, minimum width=0.465cm] at (6.177, 2.808){};
	\node[circ]  at (6.177, 2.808){};
	\node[shape=circle, xscale=-1, yscale=-1, draw, line width=1pt, minimum width=0.465cm] at (6.177, 2.058){};
	\node[circ]  at (6.177, 2.058){};
	\node[shape=circle, xscale=-1, yscale=-1, draw, line width=1pt, minimum width=0.465cm] at (4.575, 3.177){};
	\draw[line width=1pt] (4.399, 3.354) -- (4.752, 3);
	\draw[line width=1pt] (4.399, 3) -- (4.752, 3.354);
	\node[shape=circle, xscale=-1, yscale=-1, draw, line width=1pt, minimum width=0.465cm] at (4.575, 2.427){};
	\draw[line width=1pt] (4.399, 2.604) -- (4.752, 2.25);
	\draw[line width=1pt] (4.399, 2.25) -- (4.752, 2.604);
	\node[shape=circle, xscale=-1, yscale=-1, draw, line width=1pt, minimum width=0.465cm] at (6.177, 3.558){};
	\node[circ]  at (6.177, 3.558){};
	\draw (-1.25, 4) to[sinusoidal voltage source] (-1.25, 1.75);
	\draw (-1.25, 4) -| (0.802, 3.985);
	\draw (-1.25, 1.75) -| (0.802, 1.735); % node[anchor=south] at ([yshift=-0.12cm]N1.north){$i_1$};
	\draw (7.5, 3.564) -| (6.177, 3.56); %\node[flowarrow](N1) at (-0.25, 4.25){} node[anchor=south] at ([yshift=-0.12cm]N1.north){$i_1$};
	\draw (7.5, 2.048) -| (6.221, 2.062);
	\draw[line width=1pt, -latex] (7.25, 3.266) -- (7.25, 2.266) node[midway, right=2pt] {$\underline{U}_2$};
	\draw[line width=1pt, -latex] (0, 3.5) -- (0, 2.25) node[midway, left=2pt] {$\underline{U}_1$};
	\node[flowarrow](N1) at (-0.5, 4){} node[anchor=south] at ([yshift=-0.12cm]N1.north){$i_1$};
	\node[flowarrow, xscale=-1](N2) at (7.154, 3.567){} node[anchor=south] at ([yshift=-0.12cm]N2.north){$i_2$};
\end{tikzpicture}
\end{center}

Damit folgt unmittelbar für den Leerlauf ($i_2 = 0\,A$) eine
dem Windungsverhältnis proportional induzierte Spannung $\underline{U}_2$
sowie für den Kurzschlussfall ein dem Windungsverhältnis antiproportionaler
Strom.

Die Gleichungen beschreiben den Zusammenhang zwischen
Wicklungsstrom, magnetischer Feldstärke, Flussdichte und
magnetischem Fluss. Im realen Transformator wird der Fluss durch
den nichtlinearen Kernwerkstoff, den magnetischen Widerstand,
Streuanteile und gegebenenfalls die Sättigung beeinflusst. Für
den idealen Transformator wird angenommen, dass der gesamte
magnetische Fluss beide Wicklungen vollständig verkettet.

\subsection{Herleitung des Übersetzungsverhältnisses}

Primär- und Sekundärwicklung sind auf demselben magnetischen
Kreis angeordnet. Im idealen Modell verkettet deshalb derselbe
zeitabhängige Fluss $\Phi(t)$ beide Wicklungen. Wendet man das
Induktionsgesetz auf beide Wicklungen an, erhält man betragsmäßig
\[
    U_1 = N_1 \left|\frac{\mathrm d\Phi}{\mathrm dt}\right|,
    \qquad
    U_2 = N_2 \left|\frac{\mathrm d\Phi}{\mathrm dt}\right|.
\]
Das Verhältnis der beiden Spannungen ergibt sich unmittelbar,
weil der gemeinsame Flussterm herausgekürzt werden kann:
\[
    \boxed{\frac{U_1}{U_2}=\frac{N_1}{N_2}=a.}
\]
Das Übersetzungsverhältnis $a$ ist hier als Verhältnis der
Primärspannung zur Sekundärspannung beziehungsweise der
Primärwindungszahl zur Sekundärwindungszahl definiert. Bei dieser
Konvention gilt $a>1$ für einen abwärts und $a<1$ für einen
aufwärts übersetzenden Transformator.

Bei sinusförmigem magnetischem Fluss
\[
    \Phi(t)=\Phi_{\max}\sin(\omega t)
\]
ist die induzierte Spannung ebenfalls sinusförmig. Aus der
Ableitung des Flusses folgt für den Effektivwert
\[
    \boxed{U=4{,}44\,f\,N\,\Phi_{\max}.}
\]
Die Beziehung gilt für sinusförmigen Fluss und vernachlässigte
Spannungsabfälle an den Wicklungswiderständen und Streuinduktivitäten.

\subsection{Stromübersetzung und Leistungsübertragung}

Für den idealen Transformator werden Wicklungswiderstände,
Streuinduktivitäten und Kernverluste vernachlässigt. Es wird
daher keine Wirkleistung in Wärme umgesetzt. Im sinusförmigen
stationären Betrieb gilt für die Beträge der Wirkleistung bei
idealer Leistungsübertragung
\[
    U_1 I_1 = U_2 I_2.
\]
Zusammen mit dem Spannungsübersetzungsverhältnis folgt
\[
    \boxed{\frac{I_1}{I_2}=\frac{U_2}{U_1}=\frac{1}{a}.}
\]
Die Stromübersetzung verhält sich somit umgekehrt zur
Spannungsübersetzung: Eine höhere Sekundärspannung geht mit
einem entsprechend geringeren Sekundärstrom einher, sofern die
übertragene Leistung gleich bleibt.

Die Strombeziehung ist eine Betragsbeziehung. Die Vorzeichen
beziehungsweise Phasenlagen der Ströme hängen von den gewählten
Bezugsrichtungen und der Punktkennzeichnung der Wicklungen ab.


Die induzierte Spannung einer Wicklung folgt aus dem Induktionsgesetz von
Faraday, in differenzieller Form der Maxwellgleichung:

\[
    %\boxed{
        \nabla \times \vec{E}
        = -\frac{\partial \vec{B}}{\partial t}
    %}
\]

durch Integration über eine Windung und Summation über $N$ Windungen zu

\[
  u(t)=N\frac{\mathrm d\Phi(t)}{\mathrm dt},
\]
wobei das Vorzeichen von der gewählten Wicklungs- und Flussrichtung abhängt.

Bei sinusförmigem magnetischem Fluss
\[
  \Phi(t)=\Phi_{\max}\sin(\omega t)
\]
ergibt sich für den Effektivwert der induzierten Spannung
\[
  U=4{,}44\,f\,N\,\Phi_{\max}.
\]

Für zwei ideal gekoppelte Wicklungen folgt daraus
\[
  \frac{U_1}{U_2}=\frac{N_1}{N_2}=a.
\]
Das Übersetzungsverhältnis \(a\) ist hier als Verhältnis der Spannungen
auf der Seite 1 zur Seite 2 definiert.

Bei idealer Leistungsübertragung gilt
\[
  U_1 I_1=U_2 I_2
\]
für die Beträge der entsprechenden Effektivwerte. Damit folgt
\[
  \frac{I_1}{I_2}=\frac{1}{a}.
\]

Bei einem realen Drehstromtransformator müssen die Leiter-Leiter-Spannungen,
Strangspannungen, Wicklungsverschaltungen und gegebenenfalls die Phasenverschiebung
zwischen den Seiten getrennt betrachtet werden. Das einfache Verhältnis \(a\)
ist unmittelbar auf die Spannungen der jeweils gekoppelten Wicklungen anwendbar.

\subsection{Nennströme}

Für einen symmetrischen Drehstromtransformator mit Bemessungsscheinleistung
\(S_N\) und Leiter-Leiter-Bemessungsspannung \(U_N\) beträgt der Bemessungsstrom
auf der betreffenden Seite
\[
  I_N=\frac{S_N}{\sqrt{3}\,U_N}.
\]

Beispiel: Ein Transformator mit
\[
  S_N=40\,\mathrm{MVA},\qquad U_{N,1}=110\,\mathrm{kV}
\]
besitzt auf der 110-kV-Seite den Bemessungsstrom
\[
  I_{N,1}
  =\frac{40\,\mathrm{MVA}}{\sqrt{3}\cdot110\,\mathrm{kV}}
  \approx210\,\mathrm{A}.
\]

Bei \(U_{N,2}=20\,\mathrm{kV}\) ergibt sich sekundärseitig
\[
  I_{N,2}
  =\frac{40\,\mathrm{MVA}}{\sqrt{3}\cdot20\,\mathrm{kV}}
  \approx1{,}155\,\mathrm{kA}.
\]

Die Bemessungsleistung bleibt näherungsweise gleich, während sich die
Bemessungsströme umgekehrt zum Spannungsverhältnis verhalten.

\section{Reales Ersatzschaltbild}

\subsection{Serienzweig}

Das reale Ersatzschaltbild berücksichtigt die Wicklungswiderstände
\(R_1\) und \(R_2\) sowie die Streureaktanzen \(X_{\sigma 1}\) und
\(X_{\sigma 2}\).

Die Wicklungswiderstände beschreiben die ohmschen Verluste der Leiter.
Die Streureaktanzen repräsentieren den Anteil des magnetischen Flusses,
der nicht beide Wicklungen in gleicher Weise verkettet.

Für die Sekundärimpedanz gilt beim Bezug auf die Primärseite
\[
  Z_2'=a^2 Z_2.
\]
Daraus folgen für den zusammengefassten Serienzweig näherungsweise
\[
  R_k=R_1+a^2R_2,
  \qquad
  X_k=X_{\sigma 1}+a^2X_{\sigma 2},
\]
und
\[
  Z_k=R_k+jX_k.
\]

Diese Beziehungen setzen voraus, dass alle Größen auf dieselbe Seite
bezogen sind. Die Streuinduktivitäten sind dabei durch ihre Reaktanzen
bei der betrachteten Frequenz repräsentiert:
\[
  X_\sigma=\omega L_\sigma.
\]

\subsection{Magnetisierungszweig}

Der Magnetisierungszweig wird durch einen Eisenverlustwiderstand
\(R_{\mathrm{Fe}}\) und eine Magnetisierungsreaktanz \(X_m\) modelliert.
Der Widerstand repräsentiert die im Kern auftretenden Verluste,
insbesondere Hysterese- und Wirbelstromverluste. Die Reaktanz bildet
den zur Erzeugung des magnetischen Flusses erforderlichen Strom ab.

Das vereinfachte Ersatzschaltbild enthält somit:
\begin{itemize}
  \item den idealen Transformator;
  \item den Serienwiderstand und die Streureaktanz;
  \item den parallel liegenden Magnetisierungszweig.
\end{itemize}

Für Kurzschlussrechnungen wird der Magnetisierungszweig häufig vernachlässigt,
weil der Kurzschlussstrom hauptsächlich durch die Serienimpedanz begrenzt wird.
Für Leerlauf, Einschaltvorgänge und genauere Betriebsberechnungen ist diese
Vereinfachung nicht ohne Weiteres zulässig.

\section{Leerlaufversuch}

\subsection{Messprinzip}

Beim Leerlaufversuch wird eine Wicklung mit einer definierten Spannung gespeist,
während die andere Wicklung offen bleibt. Gemessen werden beispielsweise
\[
  U_0,\qquad I_0,\qquad P_0.
\]

Bei Nennfrequenz und geeigneter Prüfspannung ist der Leerlaufstrom typischerweise
deutlich kleiner als der Bemessungsstrom. Deshalb sind die ohmschen Verluste
der gespeisten Wicklung häufig klein gegenüber den Eisenverlusten.

Unter dieser Näherung gilt
\[
  P_{\mathrm{Fe}}\approx P_0.
\]

Die Näherung wird schlechter, wenn die Wicklungsverluste des Leerlaufstroms
nicht mehr vernachlässigbar sind oder die Prüfbedingungen von den
Bemessungsbedingungen abweichen.

\subsection{Bestimmung der Ersatzparameter}

Für einen einphasigen Ersatzschaltkreis ergibt sich näherungsweise
\[
  R_{\mathrm{Fe}}=\frac{U_0^2}{P_0}.
\]

Der Wirkstromanteil des Leerlaufstroms ist
\[
  I_{\mathrm{Fe}}=\frac{P_0}{U_0}.
\]

Der Magnetisierungsstrom wird bei angenommener Orthogonalität von Wirk- und
Blindstromanteil bestimmt durch
\[
  I_m=\sqrt{I_0^2-I_{\mathrm{Fe}}^2}.
\]
Damit folgt
\[
  X_m=\frac{U_0}{I_m}.
\]

Bei Drehstromtransformatoren ist zu beachten, ob die Messwerte als
Leitergrößen oder als Stranggrößen vorliegen. Die Parameterformeln sind
auf den jeweils verwendeten einphasigen Ersatzkreis anzuwenden.

\section{Kurzschlussversuch}

\subsection{Messprinzip}

Beim Kurzschlussversuch wird eine Wicklung kurzgeschlossen. An die andere
Wicklung wird eine reduzierte Spannung angelegt, sodass der gewünschte
Prüfstrom fließt. Für den üblichen Versuch mit Bemessungsstrom werden
Spannung, Strom und Wirkleistung gemessen:
\[
  U_k,\qquad I_k,\qquad P_k.
\]

Da die angelegte Spannung deutlich unter der Bemessungsspannung liegt,
ist der magnetische Fluss geringer als im Nennbetrieb. Die Eisenverluste
sind deshalb meist klein gegenüber den Wicklungsverlusten.

Die gemessene Kurzschlussleistung kann daher näherungsweise zur Bestimmung
der Serienverluste verwendet werden.

\subsection{Auswertung für einen einphasigen Ersatzkreis}

Für den einphasigen Ersatzkreis gelten
\[
  |Z_k|=\frac{U_k}{I_k},
\]
\[
  R_k=\frac{P_k}{I_k^2},
\]
und
\[
  X_k=\sqrt{|Z_k|^2-R_k^2}.
\]

Die letzte Gleichung setzt voraus, dass \(P_k\) die Wirkleistung des
betrachteten einphasigen Ersatzkreises beschreibt. Bei Drehstrommessungen
muss die Gesamtleistung konsistent auf den Strang bezogen werden.

Die Kurzschlussimpedanz lautet anschließend
\[
  Z_k=R_k+jX_k.
\]

Eine rein betragsmäßige Angabe der Kurzschlussimpedanz reicht für viele
Netzberechnungen als Näherung aus. Für genauere Berechnungen, insbesondere
bei niedrigen Spannungen oder hohem Widerstandsanteil, ist die Trennung
zwischen \(R_k\) und \(X_k\) wichtig.

\section{Kurzschlussspannung und Impedanz aus Typenschilddaten}

Die relative Kurzschlussspannung wird häufig in Prozent angegeben:
\[
  u_k=\frac{U_k}{U_N}\cdot100\%.
\]

Dabei beziehen sich \(U_k\) und \(U_N\) auf dieselbe Wicklungsseite und
dieselbe Spannungsdefinition. Bei symmetrischen Drehstromtransformatoren
kann aus der relativen Kurzschlussspannung der Betrag der Strangimpedanz
näherungsweise bestimmt werden:
\[
  |Z_k|=\frac{u_k}{100}\frac{U_N^2}{S_N}.
\]

In dieser Gleichung ist \(U_N\) die Leiter-Leiter-Bemessungsspannung
und \(S_N\) die dreiphasige Bemessungsscheinleistung. Die Gleichung liefert
die Impedanz des einphasigen Strangersatzes.

\subsection{Beispiel}

Gegeben sei ein Transformator mit
\[
  S_N=10\,\mathrm{MVA},\quad
  U_N=20\,\mathrm{kV},\quad
  u_k=8\%.
\]

Damit ergibt sich
\[
  |Z_k|
  =0{,}08\frac{(20\,\mathrm{kV})^2}{10\,\mathrm{MVA}}
  =3{,}2\,\Omega.
\]

Der Bemessungsstrom auf der 20-kV-Seite beträgt
\[
  I_N=\frac{10\,\mathrm{MVA}}{\sqrt{3}\cdot20\,\mathrm{kV}}
  \approx289\,\mathrm{A}.
\]

Der Betrag der Kurzschlussimpedanz ist damit bestimmt. Die Aufteilung
in Wirk- und Blindanteil ist aus \(u_k\) allein jedoch nicht möglich.
Dafür sind weitere Herstellerangaben oder Messdaten erforderlich.

\section{Stufenschalter und Spannungsregelung}

Stufenschalter verändern das Übersetzungsverhältnis des Transformators.
Bei einer Wicklungsanzapfung wird die wirksame Windungszahl und damit
die Übersetzung verändert. Je nach Bauart erfolgt die Umschaltung
spannungsfrei oder unter Last.

Ein Stufenschalter dient dazu, die Spannung auf der geregelten Seite
innerhalb des vorgesehenen Bereichs zu halten. Die Wirkung hängt von
Netzimpedanzen, Laststrom, Leistungsfaktor und der gewählten Regelstrategie ab.

Eine Spannungsänderung ist nicht mit einer Veränderung der Kurzschlussimpedanz
gleichzusetzen. Bei detaillierten Kurzschlussrechnungen müssen jedoch die
tatsächliche Stufenstellung und die für das Berechnungsverfahren geltenden
Vorgaben berücksichtigt werden.

\section{Schaltgruppen und Nullsystem}

Die Schaltgruppe beschreibt die Wicklungsverschaltung und die relative
Phasenverschiebung zwischen den Seiten. Sie ist für Lastfluss, Parallelbetrieb
und Fehlerberechnung von Bedeutung.

Bei der Analyse müssen mindestens folgende Eigenschaften unterschieden werden:
\begin{itemize}
  \item Phasenverschiebung zwischen Ober- und Unterspannungsseite;
  \item zugänglicher Neutralpunkt und dessen Erdung;
  \item Möglichkeit von Nullsystemströmen in den Wicklungen;
  \item magnetischer Rückweg des Nullsystemflusses.
\end{itemize}

Eine Dreieckwicklung kann einen Umlaufweg für bestimmte Nullsystemströme
bereitstellen. Daraus folgt aber nicht automatisch, dass Nullsystemströme
zwischen den beiden Netzseiten übertragen werden. Das Verhalten ist anhand
der konkreten Schaltgruppe, der Wicklungsanordnung und der Erdung zu beurteilen.

\section{Parallelbetrieb}

Für den Parallelbetrieb von Transformatoren sind insbesondere zu prüfen:
\begin{itemize}
  \item gleiches oder kompatibles Übersetzungsverhältnis;
  \item passende Phasenverschiebung und Schaltgruppe;
  \item gleiche Phasenfolge und richtige Klemmenzuordnung;
  \item vergleichbare Kurzschlussspannungen;
  \item geeignete Impedanzwinkel und Bemessungsleistungen.
\end{itemize}

Abweichende Übersetzungsverhältnisse können Ausgleichsströme verursachen.
Abweichende Kurzschlussimpedanzen führen zu einer ungleichen Aufteilung
der Lastströme. Ein Parallelbetrieb darf daher nicht allein aufgrund
gleicher Nennspannungen angenommen werden.

\section{Zusammenfassung}

Für Netzberechnungen ist die Transformatorimpedanz ein zentrales Element.
Die Kurzschlussspannung liefert deren Betrag auf der gewählten Bezugsseite.
Die Trennung in Wirk- und Blindanteil erfordert zusätzliche Informationen.
Schaltgruppe und Erdung sind besonders für Nullsystem- und Erdschlussmodelle
entscheidend.

\section{Übungsaufgaben}

\begin{enumerate}
  \item Berechnen Sie die Bemessungsströme eines Transformators mit
        \(S_N=25\,\mathrm{MVA}\) und \(U_N=110/20\,\mathrm{kV}\).
  \item Ein Transformator besitzt \(S_N=16\,\mathrm{MVA}\),
        \(U_N=10\,\mathrm{kV}\) und \(u_k=6\%\).
        Bestimmen Sie den Betrag der Strang-Kurzschlussimpedanz auf der
        10-kV-Seite.
  \item Erläutern Sie, warum sich aus \(u_k\) allein nicht bestimmen lässt,
        welcher Anteil der Kurzschlussimpedanz ohmsch ist.
  \item Beschreiben Sie, welche Informationen zur Bewertung des
        Nullsystemverhaltens eines Transformators erforderlich sind.
\end{enumerate}

\chapter{Leitungen und Kabel}
\label{chap:05-leitungen-und-kabel}

\section{Einordnung und Modellierungsziele}

Leitungen und Kabel verbinden Betriebsmittel und bestimmen durch ihre
Serienimpedanzen und Queradmittanzen den Spannungsfall, die Verluste,
die Kurzschlussstromhöhe und das Erdschlussverhalten.

Freileitungen und Kabel unterscheiden sich unter anderem in ihrer Geometrie,
ihrem Isolationssystem, ihrer Wärmeabfuhr und ihrer kapazitiven Kopplung
zur Erde. Die Modellierung muss deshalb zum betrachteten Phänomen passen.

Für eine erste Netzberechnung wird häufig ein einphasiges Ersatzmodell
je Phase verwendet. Bei unsymmetrischen Fehlern und Erdschlussberechnungen
müssen zusätzlich die Sequenzimpedanzen und die Rückstrompfade berücksichtigt werden.

\section{Längs- und Querparameter}

Die längenbezogenen Serienparameter einer Leitung werden durch
\[
  Z'=R'+j\omega L'
\]
beschrieben. Die Queradmittanz lautet
\[
  Y'=G'+j\omega C'.
\]

Dabei bezeichnen:
\begin{itemize}
  \item \(R'\): längenbezogener Wirkwiderstand;
  \item \(L'\): längenbezogene Induktivität;
  \item \(G'\): längenbezogener Ableitwert;
  \item \(C'\): längenbezogene Kapazität;
  \item \(\omega=2\pi f\): Kreisfrequenz.
\end{itemize}

Für eine Leitung der Länge \(l\) ergibt das konzentrierte Modell
\[
  Z=Z'l,\qquad Y=Y'l.
\]

Diese Gleichungen setzen voraus, dass die Parameter entlang der Leitung
hinreichend gleichmäßig sind. Bei langen Leitungen oder wenn die Queradmittanz
wesentlich ist, ist eine verteilte Leitungsbeschreibung genauer.

\section{Ohmscher Widerstand}

\subsection{Grundgleichung}

Für einen homogenen Leiter mit Länge \(l\), Querschnitt \(A\) und
spezifischem Widerstand \(\rho\) gilt
\[
  R=\rho\frac{l}{A}.
\]

Der Widerstand wächst proportional mit der Länge und sinkt bei größerem
Querschnitt. Die Materialeigenschaften gehen über \(\rho\) ein.

Bei Wechselstrom können Skin- und Proximity-Effekte den wirksamen Widerstand
gegenüber dem Gleichstromwert erhöhen. Bei mehradrigen Kabeln und dicht
benachbarten Leitern ist deshalb nicht immer der reine Gleichstromwiderstand
für Verlust- und Spannungsfallberechnungen ausreichend.

\subsection{Temperaturabhängigkeit}

In einem begrenzten Temperaturbereich lässt sich der Widerstand näherungsweise
mit
\[
  R_\vartheta=R_{\vartheta_0}
  \left[1+\alpha(\vartheta-\vartheta_0)\right]
\]
korrigieren. \(\alpha\) ist der materialabhängige Temperaturkoeffizient.

Die tatsächliche Leitertemperatur ist nicht nur eine Eingangsgröße:
Sie stellt sich aus Verlustleistung, thermischem Widerstand der Umgebung,
Verlegebedingungen und Wärmeabfuhr ein. Für die Bemessung ist deshalb eine
thermische Betrachtung erforderlich.

\section{Induktivität und Reaktanz}

Die Induktivität hängt von Leitergeometrie, Leiterabständen, Rückstrompfaden
und magnetischen Eigenschaften der Umgebung ab. Für die Serienreaktanz gilt
\[
  X_L=\omega L.
\]

Die Reaktanz wächst bei gleichbleibender Induktivität linear mit der Frequenz.
Bei Wechselstromnetzen beeinflusst sie den Spannungsfall und den Blindleistungsfluss.

Bei symmetrischen Drehstromleitungen wird die gegenseitige magnetische Kopplung
der Phasen in den Ersatzparametern berücksichtigt. Bei unsymmetrischer Geometrie
oder asymmetrischer Leitungsführung kann eine einfache Einphasenbetrachtung
ungenau sein.

\section{Kapazität und Ableitwert}

\subsection{Kapazitive Kopplung}

Zwischen Leitern sowie zwischen Leitern und geerdeten Teilen entstehen
elektrische Felder. Daraus folgt eine Kapazität. Der kapazitive Strom eines
einfachen Zweiges beträgt
\[
  I_C=\omega C U.
\]

Die verwendete Spannung \(U\) muss zur Definition der Kapazität passen.
Eine Kapazität zwischen zwei Phasen wird nicht ohne Umrechnung wie eine
Kapazität einer Phase gegen Erde behandelt.

Der Ableitwert \(G\) beschreibt reale dielektrische Verluste und Leckpfade.
Bei vielen vereinfachten Netzberechnungen ist er gegenüber dem kapazitiven
Anteil klein, kann bei speziellen Isolations- und Betriebsbedingungen aber
relevant sein.

\subsection{Kabel gegenüber Freileitungen}

Kabel besitzen aufgrund der geringen Leiterabstände und der hohen relativen
Permittivität ihrer Isolierung typischerweise eine deutlich größere Kapazität
je Längeneinheit als Freileitungen.

Damit steigt bei gleicher Spannung und Frequenz der kapazitive Ladestrom
mit der Kabellänge. In Mittelspannungsnetzen kann dieser Zusammenhang das
Erdschlussverhalten wesentlich beeinflussen.

\section{Leitungsmodelle}

\subsection{Kurze Leitung}

Bei einer kurzen Leitung kann die Queradmittanz für bestimmte Fragestellungen
vernachlässigt werden. Das Ersatzmodell besteht dann aus einer Serienimpedanz
\[
  Z=R+jX.
\]

Dieses Modell ist besonders für überschlägige Spannungsfall- und
Kurzschlussberechnungen nützlich. Es bildet jedoch keine verteilten
kapazitiven Ströme ab.

\subsection{Nominales Pi-Modell}

Bei einer Leitung mittlerer Länge wird die Queradmittanz häufig auf beide
Leitungsenden aufgeteilt. Das nominale Pi-Modell enthält die Serienimpedanz
\(Z\) und zwei Querzweige mit jeweils \(Y/2\).

Die Näherung bildet den Einfluss der Leitungskapazität besser ab als das
reine Serienmodell. Sie ist jedoch nicht identisch mit einer vollständigen
verteilten Leitungsrechnung.

\subsection{Verteilte Parameter}

Bei einer Leitung mit gleichmäßig verteilten Parametern werden Spannung
und Strom entlang der Leitung durch die Leitungsgleichungen beschrieben.
Mit den längenbezogenen Größen \(Z'\) und \(Y'\) ist die Ausbreitungskonstante
\[
  \gamma=\sqrt{Z'Y'}
\]
und der Wellenwiderstand
\[
  Z_W=\sqrt{\frac{Z'}{Y'}}.
\]

Die Zweitorbeziehung für eine homogene Leitung der Länge \(l\) kann in der Form
\[
  \begin{bmatrix}
    U_1\\ I_1
  \end{bmatrix}
  =
  \begin{bmatrix}
    \cosh(\gamma l) & Z_W\sinh(\gamma l)\\
    \frac{1}{Z_W}\sinh(\gamma l) & \cosh(\gamma l)
  \end{bmatrix}
  \begin{bmatrix}
    U_2\\ I_2
  \end{bmatrix}
\]
geschrieben werden, sofern die Stromrichtungen und die Zuordnung der
Endgrößen konsistent definiert sind.

Dieses Modell ist insbesondere für lange Leitungen relevant. Die genaue
Anwendung erfordert eine einheitliche Vorzeichenkonvention für die Ströme.

\section{Spannungsfall bei Drehstrom}

Für eine kurze, symmetrisch belastete Drehstromleitung kann der Spannungsfall
näherungsweise durch
\[
  \Delta U\approx\sqrt{3}\,I\,l
  \left(R'\cos\varphi+X'\sin\varphi\right)
\]
abgeschätzt werden.

Hier sind \(R'\) und \(X'\) die längenbezogenen Serienparameter je Phase,
\(I\) der Leiterstrom und \(l\) die einfache Leitungslänge. Die Formel
verwendet die übliche Näherung für kleine Spannungsänderungen und einen
definierten Leistungsfaktorwinkel \(\varphi\).

Bei induktiver Last ist \(\sin\varphi\) positiv. Bei führendem Leistungsfaktor
kann der Blindleistungsanteil das Vorzeichen wechseln. Eine kapazitive
Einspeisung kann daher den Spannungsfall verringern oder eine Spannungsanhebung
verursachen.

Für größere Spannungsänderungen, lange Leitungen oder bedeutende Queradmittanzen
ist eine Phasorrechnung oder ein vollständiges Leitungsmodell zu verwenden.

\section{Leitungsverluste}

Die ohmschen Verluste eines symmetrischen Drehstromsystems sind näherungsweise
\[
  P_{\mathrm V}=3I^2R.
\]

Dabei ist \(R\) der Widerstand eines Leiters über die betrachtete Leitungslänge.
Zusätzliche Verluste können in metallischen Kabelschirmen, Bewehrungen,
Erdungsleitern und anderen leitfähigen Komponenten auftreten.

Für eine gegebene übertragene Wirkleistung \(P\) und einen gegebenen
Leistungsfaktor gilt
\[
  I=\frac{P}{\sqrt{3}\,U\,\cos\varphi}.
\]
Ein höheres Spannungsniveau reduziert daher bei gleicher Leistung und gleichem
Leistungsfaktor den Strom und damit die ohmschen Verluste, sofern die übrigen
Parameter vergleichbar bleiben.

\section{Thermische Strombelastbarkeit}

Die zulässige Dauerstrombelastbarkeit ist eine thermische Bemessungsgröße.
Sie hängt insbesondere ab von:
\begin{itemize}
  \item Leiterwerkstoff und Leiterquerschnitt;
  \item zulässiger Leitertemperatur und Isolationssystem;
  \item Umgebungstemperatur und Verlegeart;
  \item Häufung mehrerer belasteter Kabel;
  \item thermischem Widerstand des Bodens bei Erdkabeln;
  \item Verlusten in Schirmen und metallischen Umhüllungen;
  \item Lastprofil und thermischen Randbedingungen.
\end{itemize}

Eine Querschnittswahl allein anhand des ohmschen Widerstands ist deshalb
nicht ausreichend. Zusätzlich sind Spannungsfall, Kurzschlussfestigkeit,
Schutzbedingungen, mechanische Anforderungen und die einschlägigen
Bemessungsregeln zu berücksichtigen.

\section{Kurzschlussbeanspruchung von Leitern}

Während eines Kurzschlusses kann der Strom wesentlich höher als der
Dauerstrom sein. Für eine vereinfachte adiabatische Betrachtung wird häufig
eine Beziehung der Form
\[
  S=\frac{I_k\sqrt{t}}{k}
\]
verwendet, wobei \(S\) den Leiterquerschnitt, \(I_k\) den wirksamen
Kurzschlussstrom, \(t\) die Fehlerdauer und \(k\) einen vom Werkstoff,
Isolationssystem und den zulässigen Anfangs- und Endtemperaturen abhängigen
Faktor bezeichnet.

Diese Beziehung darf nur mit einem zum Leiter und zur Isolierung passenden
Faktor angewendet werden. Bei stark zeitabhängigen Strömen, mehreren
Fehlerbeiträgen oder nichtadiabatischen Bedingungen sind die jeweils
zutreffenden Verfahren und Normen heranzuziehen.

\section{Kabelschirme und Erdung}

Metallische Schirme und Bewehrungen erfüllen elektrische, mechanische und
feldsteuernde Aufgaben. Ihre Erdung und Verbindung beeinflussen:
\begin{itemize}
  \item die Sicherheit bei Isolationsfehlern;
  \item die Höhe und Verteilung von Schirmströmen;
  \item die induzierten Spannungen;
  \item die Nullsystemimpedanz;
  \item die thermische Beanspruchung bei Erd- und Kurzschlüssen.
\end{itemize}

Beidseitige Schirmerdung kann unter bestimmten Bedingungen geschlossene
Strompfade für induzierte Ströme erzeugen. Einseitige Erdung reduziert solche
Schirmströme möglicherweise, kann aber am nicht geerdeten Ende höhere
induzierte Spannungen verursachen. Die Wahl der Schirmerdung ist deshalb
eine Systementscheidung und keine allgemeine Regel, die unabhängig von
Kabellänge, Spannungsebene und Fehlerfall getroffen werden kann.

\section{Übungsaufgaben}

\begin{enumerate}
  \item Berechnen Sie den Gleichstromwiderstand eines Kupferleiters mit
        vorgegebenem Querschnitt und einer Länge von \(250\,\mathrm m\).
        Berücksichtigen Sie anschließend eine Temperaturerhöhung.
  \item Ein Kabel besitzt eine Kapazität von \(0{,}25\,\mathrm{\mu F/km}\)
        je Phase gegen Erde. Berechnen Sie den kapazitiven Strom für eine
        Leitungslänge von \(8\,\mathrm{km}\) bei \(20\,\mathrm{kV}\) und
        \(50\,\mathrm{Hz}\). Geben Sie an, welche Spannung Sie einsetzen.
  \item Erläutern Sie die Unterschiede zwischen Serienmodell, nominalem
        Pi-Modell und verteiltem Leitungsmodell.
  \item Warum kann die Schirmerdung die Erdschlussberechnung beeinflussen?
  \item Welche Informationen fehlen, wenn ein Kabel allein anhand des
        Querschnitts thermisch bemessen werden soll?
\end{enumerate}
\chapter{Ersatzschaltbilder der Netztechnik}
\label{chap:06-ersatzschaltbilder}

%\chapter{Zusammengesetzte Ersatzschaltbilder und Netzreduktion}
%\label{chap:ersatznetze}

\section{Ziel der Netzmodellierung}

Ein elektrisches Netz besteht aus Quellen, Transformatoren, Leitungen,
Kabeln, Lasten und gegebenenfalls rotierenden Maschinen oder Umrichtern.
Für eine konkrete Berechnung wird daraus ein Modell gebildet, das die für
die jeweilige Fragestellung relevanten Eigenschaften erhält.

Das Modell eines Kurzschlusses muss nicht dieselben Details enthalten wie
ein Modell für Spannungsfall, Verlustleistung oder Erdschlussstrom. Die
Modellierung beginnt deshalb mit der Frage, welche Größe bestimmt werden soll.

\section{Impedanzen auf eine gemeinsame Spannungsebene beziehen}

Für einen idealen Transformator mit Übersetzungsverhältnis
\[
  a=\frac{U_1}{U_2}
\]
wird eine sekundärseitige Impedanz auf die Primärseite bezogen durch
\[
  Z_2'=a^2Z_2.
\]
Umgekehrt wird eine primärseitige Impedanz auf die Sekundärseite bezogen durch
\[
  Z_1'=\frac{Z_1}{a^2}.
\]

Diese Regeln gelten für die Impedanzen der jeweils gekoppelten Ersatzkreise.
Bei Drehstromtransformatoren sind die Spannungs- und Strangdefinitionen
konsistent zu halten.

\subsection{Beispiel}

Eine Impedanz auf der 0{,}4-kV-Seite eines idealisierten Transformators
soll auf die 20-kV-Seite bezogen werden. Mit dem Verhältnis
\[
  a=\frac{20}{0{,}4}=50
\]
gilt
\[
  Z'=50^2 Z=2500Z.
\]

Eine sekundärseitige Impedanz von \(0{,}01\,\Omega\) entspricht damit
auf der 20-kV-Seite einer Impedanz von
\[
  Z'=25\,\Omega.
\]

Die Umrechnung verändert nicht das physikalische Verhalten des Netzes.
Sie verändert nur die Darstellung der Impedanz auf einer anderen
Spannungsseite.

\section{Komplexe Impedanzen addieren}

Für eine Reihenschaltung von Impedanzen gilt
\[
  Z_{\mathrm{ges}}=\sum_{i=1}^{n}Z_i.
\]

Mit
\[
  Z_i=R_i+jX_i
\]
folgt
\[
  R_{\mathrm{ges}}=\sum_i R_i,\qquad
  X_{\mathrm{ges}}=\sum_i X_i.
\]
Erst anschließend wird der Betrag bestimmt:
\[
  |Z_{\mathrm{ges}}|=
  \sqrt{R_{\mathrm{ges}}^2+X_{\mathrm{ges}}^2}.
\]

Im Allgemeinen gilt nicht
\[
  |Z_1+Z_2|=|Z_1|+|Z_2|.
\]
Die Beträge einzelner komplexer Impedanzen dürfen deshalb nicht einfach
addiert werden.

Für eine Parallelschaltung gilt
\[
  \frac{1}{Z_{\mathrm{ges}}}=\sum_i\frac{1}{Z_i}.
\]
Es ist häufig zweckmäßig, zuerst die Admittanzen
\[
  Y_i=\frac{1}{Z_i}
\]
zu bilden und diese zu addieren.

\section{Thévenin-Ersatzquelle}

Ein lineares Netzwerk kann aus Sicht eines betrachteten Anschlusses unter
geeigneten Voraussetzungen durch eine Ersatzspannungsquelle und eine
Serienimpedanz dargestellt werden.

Die Thévenin-Ersatzspannung \(U_{\mathrm{Th}}\) ist die Leerlaufspannung
am betrachteten Anschluss. Die Thévenin-Impedanz \(Z_{\mathrm{Th}}\)
ist die von diesem Anschluss aus gesehene Impedanz bei deaktivierten
unabhängigen Quellen. Ideale Spannungsquellen werden dabei kurzgeschlossen,
ideale Stromquellen geöffnet. Abhängige Quellen müssen bei der Bestimmung
der Ersatzimpedanz weiterhin berücksichtigt werden.

Bei der Kurzschlussberechnung wird die Ersatzimpedanz aus Sicht der
Fehlerstelle gebildet. Dadurch lässt sich ein komplexes Netz auf ein
einfaches Ersatzmodell reduzieren.

\section{Kurzschlussstrom im Ersatznetz}

Für einen einphasigen Ersatzkreis mit Fehlerimpedanz \(Z_f\) gilt
\[
  I_k=\frac{U_{\mathrm{Th}}}{Z_{\mathrm{Th}}+Z_f}.
\]

Bei einem symmetrischen dreipoligen Fehler wird das Mitsystem-Ersatznetz
verwendet. Die Spannung ist dann die Strangspannung:
\[
  U_{\mathrm{ph}}=\frac{U_{LL}}{\sqrt{3}}.
\]
Daher gilt für den Betrag des symmetrischen Kurzschlussstroms
\[
  I_k=\frac{U_{LL}}{\sqrt{3}|Z_{\mathrm{Th}}+Z_f|}.
\]

Diese Gleichung beschreibt das vereinfachte Ersatznetz. Bei normgerechten
Kurzschlussberechnungen können je nach Berechnungsziel zusätzliche Faktoren,
Anfangsbedingungen, Motorbeiträge, Spannungsfaktoren und Korrekturen
notwendig sein.

\section{Per-Unit-System als einheitliche Darstellung}

Mit einer gewählten dreiphasigen Bezugsleistung \(S_B\) und einer
Leiter-Leiter-Bezugsspannung \(U_B\) werden die Bezugsgrößen definiert als
\[
  I_B=\frac{S_B}{\sqrt{3}U_B},
  \qquad
  Z_B=\frac{U_B^2}{S_B}.
\]

Der Per-Unit-Wert einer Impedanz ist
\[
  z_{\mathrm{pu}}=\frac{Z}{Z_B}.
\]

Für einen Wechsel der Bezugsbasis gilt
\[
  z_{\mathrm{pu,new}}
  =z_{\mathrm{pu,old}}
   \frac{S_{B,\mathrm{new}}}{S_{B,\mathrm{old}}}
   \left(\frac{U_{B,\mathrm{old}}}{U_{B,\mathrm{new}}}\right)^2.
\]

Werden die Bezugsspannungen über ideale Transformatoren konsistent gewählt,
bleiben Per-Unit-Impedanzen auf beiden Seiten häufig besonders einfach
vergleichbar. Die Bezugsspannungen müssen dabei dem Übersetzungsverhältnis
entsprechen.

\section{Durchgängiges Rechenbeispiel}

Gegeben sei ein vereinfachtes Netz mit:
\begin{itemize}
  \item Einspeisung mit Kurzschlussleistung \(S_k''=250\,\mathrm{MVA}\)
        auf der 20-kV-Seite;
  \item Transformator \(20/0{,}4\,\mathrm{kV}\);
  \item Kabel auf der 0{,}4-kV-Seite;
  \item dreipoligem Kurzschluss am Kabelende.
\end{itemize}

Die Netzimpedanz wird zunächst auf der 20-kV-Seite abgeschätzt:
\[
  |Z_{\mathrm{Netz}}|
  \approx\frac{U^2}{S_k''}
  =\frac{(20\,\mathrm{kV})^2}{250\,\mathrm{MVA}}
  =1{,}6\,\Omega.
\]

Diese Formel liefert eine Näherung für den Betrag der Netzimpedanz bei
der betrachteten Spannungsebene. Sie bestimmt nicht deren Wirk- und
Blindanteil getrennt.

Für die Übertragung auf die 0{,}4-kV-Seite ist das Spannungsverhältnis
\[
  a=\frac{20}{0{,}4}=50.
\]
Damit ergibt sich näherungsweise
\[
  |Z_{\mathrm{Netz}}'|=\frac{1{,}6}{50^2}
  =0{,}00064\,\Omega.
\]

Anschließend werden die auf die Niederspannungsseite bezogene Netzimpedanz,
die Transformatorimpedanz und die Kabelimpedanz komplex addiert:
\[
  Z_{\mathrm{ges}}
  =Z_{\mathrm{Netz}}'+Z_{\mathrm{Tr}}+Z_{\mathrm{Kabel}}.
\]

Der anfängliche symmetrische Kurzschlussstrom ergibt sich im vereinfachten
Ersatzmodell zu
\[
  I_k''=\frac{U_{LL}}{\sqrt{3}|Z_{\mathrm{ges}}|}.
\]

Ein numerischer Endwert ist erst dann fachlich belastbar, wenn die
komplexen Transformator- und Kabelimpedanzen sowie die verwendeten
Kurzschlussannahmen vollständig bekannt sind.

\section{Plausibilitätsprüfungen}

Nach einer Netzberechnung sollten mindestens folgende Prüfungen erfolgen:
\begin{itemize}
  \item Stimmen Spannungs-, Strom- und Impedanzdimensionen?
  \item Wurden alle Impedanzen auf dieselbe Spannungsebene bezogen?
  \item Wurden komplexe Impedanzen und nicht nur ihre Beträge addiert?
  \item Ist der errechnete Kurzschlussstrom kleiner als die durch die
        Ersatzimpedanz begrenzte plausible Größenordnung?
  \item Sind die Beiträge von Motoren, Generatoren und Umrichtern
        im verwendeten Modell berücksichtigt oder bewusst vernachlässigt?
  \item Entspricht die verwendete Spannung der Definition als Strang-
        oder Leiter-Leiter-Spannung?
\end{itemize}

\section{Übungsaufgaben}

\begin{enumerate}
  \item Beziehen Sie eine Impedanz von \(0{,}05\,\Omega\) auf der
        0{,}4-kV-Seite auf die 20-kV-Seite.
  \item Addieren Sie \(Z_1=0{,}1+j0{,}3\,\Omega\) und
        \(Z_2=0{,}2+j0{,}1\,\Omega\). Bestimmen Sie Betrag und Winkel
        der Gesamtimpedanz.
  \item Erklären Sie die Bedeutung der Thévenin-Ersatzimpedanz für die
        Berechnung eines Kurzschlussstroms.
  \item Welche zusätzlichen Informationen sind erforderlich, um das
        durchgerechnete Netzbeispiel mit einem belastbaren Zahlenwert
        abzuschließen?
\end{enumerate}

\part{Kurzschluss und Erdschluss}
\chapter{Symmetrische Komponenten}
\label{chap:07-symmetrische-komponenten}

\section{Motivation}

Die Methode der symmetrischen Komponenten ermöglicht es, ein unsymmetrisches
Dreiphasensystem in drei symmetrische Teilsysteme zu zerlegen. Sie ist ein
zentrales Verfahren zur Analyse unsymmetrischer Kurzschlüsse und Erdschlüsse.

Die Methode ist besonders leistungsfähig, wenn die Betriebsmittel und das
Netz im betrachteten Modell linear beschrieben werden können. Bei stark
nichtlinearen Betriebsmitteln oder leistungselektronischen Einspeisern muss
geprüft werden, ob das zugrunde gelegte Modell die tatsächliche Fehlerreaktion
hinreichend beschreibt.

\section{Definition der symmetrischen Komponenten}

Es sei
\[
  a=e^{j120^\circ}
  =-\frac12+j\frac{\sqrt{3}}{2}.
\]
Dann gilt
\[
  a^2=e^{j240^\circ}
  =-\frac12-j\frac{\sqrt{3}}{2}
\]
sowie
\[
  a^3=1,\qquad 1+a+a^2=0.
\]

Die Phasenspannungen \(U_a,U_b,U_c\) werden in Null-, Mit- und
Gegensystemkomponenten zerlegt:
\[
  \begin{bmatrix}
    U_0\\
    U_1\\
    U_2
  \end{bmatrix}
  =
  \frac13
  \begin{bmatrix}
    1&1&1\\
    1&a&a^2\\
    1&a^2&a
  \end{bmatrix}
  \begin{bmatrix}
    U_a\\
    U_b\\
    U_c
  \end{bmatrix}.
\]

Die Rücktransformation lautet
\[
  \begin{bmatrix}
    U_a\\
    U_b\\
    U_c
  \end{bmatrix}
  =
  \begin{bmatrix}
    1&1&1\\
    1&a^2&a\\
    1&a&a^2
  \end{bmatrix}
  \begin{bmatrix}
    U_0\\
    U_1\\
    U_2
  \end{bmatrix}.
\]

Dieselbe Transformation kann für die Ströme verwendet werden. Die
Phasenfolge und die Vorzeichenkonvention müssen dabei durchgehend
konsistent sein.

\section{Physikalische Bedeutung}

\subsection{Mitsystem}

Das Mitsystem besitzt die normale Phasenfolge \(a\)-\(b\)-\(c\).
Bei einem symmetrischen Netzbetrieb ist es das dominierende System.

\subsection{Gegensystem}

Das Gegensystem besitzt die umgekehrte Phasenfolge. Es wird insbesondere
bei unsymmetrischen Belastungen und Fehlern angeregt.

Bei rotierenden Maschinen können Gegensystemströme zusätzliche Verluste
und thermische Beanspruchungen verursachen. Die zulässige Beanspruchung
ist deshalb betriebsmittelabhängig.

\subsection{Nullsystem}

Im Nullsystem sind alle drei Phasoren gleichphasig:
\[
  U_{a0}=U_{b0}=U_{c0}=U_0.
\]

Die Nullsystemkomponente ist eng mit dem Neutralleiter- und Erdungsrückweg
verbunden. Ob ein Nullsystemstrom fließen kann und wie groß die
Nullsystemimpedanz ist, hängt von der Netz- und Betriebsmittelstruktur ab.

\section{Nullsystem und Neutralleiterstrom}

Für die Phasenströme gilt
\[
  I_a+I_b+I_c=3I_0.
\]

Der Neutralleiterstrom hängt von der gewählten Stromrichtung ab. Wird
der Neutralleiterstrom vom Sternpunkt zum Netz als positiv definiert,
kann beispielsweise
\[
  I_N=-(I_a+I_b+I_c)=-3I_0
\]
gelten.

Die Gleichung verdeutlicht, dass sich die drei Phasenströme bei einem
symmetrischen Betrieb ohne Nullsystemkomponente vektoriell zu null addieren.
Bei einem Erdschluss oder einer unsymmetrischen Belastung kann eine
Nullsystemkomponente auftreten.

\section{Sequenzimpedanzen}

Für die drei Sequenzsysteme werden getrennte Ersatzimpedanzen angesetzt:
\[
  Z_1,\qquad Z_2,\qquad Z_0.
\]

Die Mitsystemimpedanz \(Z_1\) beschreibt die Reaktion auf Mitsystemgrößen.
Die Gegensystemimpedanz \(Z_2\) beschreibt die Reaktion auf Gegensystemgrößen.
Die Nullsystemimpedanz \(Z_0\) hängt wesentlich von den Nullsystem-Rückwegen ab.

Für viele passive Betriebsmittel sind Mit- und Gegensystemimpedanzen ähnlich.
Diese Ähnlichkeit darf jedoch nicht ohne Prüfung als exakte Gleichheit
angenommen werden. Die Nullsystemimpedanz kann deutlich anders sein.

Bei Kabeln beeinflussen insbesondere metallische Schirme, Bewehrungen,
Erdungsverbindungen und die Anordnung der Leiter das Nullsystemverhalten.
Bei Transformatoren sind Wicklungsschaltung, Sternpunktzugänglichkeit,
Erdung und magnetischer Rückweg zu berücksichtigen.

\section{Sequenznetze und Fehlerbedingungen}

Die Sequenznetze werden entsprechend den Randbedingungen des betrachteten
Fehlers miteinander verbunden. Dabei werden die physikalischen
Fehlerbedingungen zunächst in Phasengrößen formuliert und anschließend
in Sequenzgrößen übersetzt.

Ein dreipoliger symmetrischer Kurzschluss beansprucht im idealisierten
symmetrischen Modell nur das Mitsystem. Ein einpoliger Erdschluss erfordert
dagegen im Allgemeinen die Mit-, Gegen- und Nullsysteme.

\section{Beispiel: Zerlegung eines unsymmetrischen Systems}

Gegeben seien die Phasenspannungen
\[
  U_a=230\angle 0^\circ\,\mathrm V,
\]
\[
  U_b=200\angle -125^\circ\,\mathrm V,
\]
\[
  U_c=240\angle 118^\circ\,\mathrm V.
\]

Die Sequenzkomponenten ergeben sich aus
\[
  U_0=\frac13(U_a+U_b+U_c),
\]
\[
  U_1=\frac13(U_a+aU_b+a^2U_c),
\]
\[
  U_2=\frac13(U_a+a^2U_b+aU_c).
\]

Da die Größen als komplexe Phasoren gegeben sind, müssen sie vor der
Addition in kartesische Form oder mit einer konsistenten komplexen
Rechenoperation überführt werden. Die Beträge allein dürfen nicht addiert
werden.

Die Rücktransformation muss anschließend wieder die ursprünglichen
Phasenspannungen ergeben. Sie ist eine nützliche Plausibilitätsprüfung
für eine manuelle oder softwaregestützte Rechnung.

\section{Grenzen der Methode}

Die symmetrischen Komponenten sind eine mathematische Transformation
der Phasengrößen. Sie machen ein unsymmetrisches Problem nicht automatisch
linear und beseitigen auch nicht die Notwendigkeit korrekter Betriebsmittelmodelle.

Besonders zu prüfen sind:
\begin{itemize}
  \item die Gültigkeit linearer Ersatzimpedanzen;
  \item die korrekte Nullsystemmodellierung;
  \item die Phasenfolge und die Vorzeichenkonvention;
  \item die Auswirkungen von Transformator-Schaltgruppen;
  \item die Fehlerimpedanz und der konkrete Erdungszustand;
  \item die Eigenschaften von Umrichtern und anderen geregelten Quellen.
\end{itemize}

\section{Übungsaufgaben}

\begin{enumerate}
  \item Leiten Sie die Beziehung \(I_a+I_b+I_c=3I_0\) aus der
        Rücktransformation her.
  \item Beschreiben Sie die physikalische Bedeutung von Mit-, Gegen-
        und Nullsystem.
  \item Warum ist die Nullsystemimpedanz eines Kabelsystems nicht
        ohne Weiteres gleich der Mitsystemimpedanz?
  \item Erläutern Sie, warum ein symmetrischer dreipoliger Kurzschluss
        im idealisierten Modell nur das Mitsystem benötigt.
\end{enumerate}
\chapter{Kurzschlussberechnung}
\label{chap:08-kurzschlussberechnung}

%\chapter{Dreipolige Kurzschlussberechnung}
%\label{chap:kurzschluss}

\section{Ziel und Einordnung}

Ein dreipoliger Kurzschluss verbindet die drei Phasen an einer Fehlerstelle.
Bei ideal symmetrischen Netzbedingungen ist der Fehler symmetrisch.
Die Berechnung kann dann mit einem einphasigen Ersatzkreis erfolgen.

Der dreipolige Kurzschluss führt häufig zu hohen Strömen und ist für die
Bemessung von Schaltgeräten, Sammelschienen, Leitern und Schutzsystemen
von Bedeutung. Er ist jedoch nicht in jedem Netz und für jede Fehlerart
der Fehler mit dem höchsten Strom.

\section{Ersatzschaltbild und Anfangskurzschlussstrom}

Das Netz wird aus Sicht der Fehlerstelle durch eine Ersatzspannung und
eine Ersatzimpedanz beschrieben. Für den symmetrischen Drehstromfall
ist die Strangspannung
\[
  U_{\mathrm{ph}}=\frac{U_{LL}}{\sqrt{3}}.
\]

Der anfängliche symmetrische Kurzschlussstrom kann in einem vereinfachten
Ersatzmodell durch
\[
  I_k''=\frac{U_{\mathrm{ph}}}{|Z_k|}
  =\frac{U_{LL}}{\sqrt{3}|Z_k|}
\]
bestimmt werden.

Das Symbol \(I_k''\) wird in der Kurzschlussberechnung häufig für den
anfänglichen symmetrischen Kurzschlussstrom verwendet. Bei normgerechten
Berechnungen sind die jeweils geltenden Definitionen und Korrekturfaktoren
zu berücksichtigen.

\section{Kurzschlussleistung}

Die dreiphasige Kurzschlussleistung wird häufig definiert durch
\[
  S_k''=\sqrt{3}\,U_{LL}I_k''.
\]

Daraus folgt
\[
  I_k''=\frac{S_k''}{\sqrt{3}\,U_{LL}}.
\]

Unter der vereinfachenden Annahme eines symmetrischen Netzes und einer
geeigneten Ersatzimpedanz kann der Betrag der Netzimpedanz abgeschätzt
werden als
\[
  |Z_k|\approx\frac{U_{LL}^2}{S_k''}.
\]

Die Kurzschlussleistung ist eine Rechengröße, die den Kurzschlussstrom
mit der Bezugsspannung verknüpft. Sie ist nicht mit der während des
gesamten Fehlerverlaufs konstant abgegebenen Wirkleistung gleichzusetzen.

\section{Beispiel: Kurzschlussleistung auf Mittelspannungsebene}

Gegeben seien
\[
  U_{LL}=20\,\mathrm{kV},\qquad S_k''=250\,\mathrm{MVA}.
\]

Der anfängliche symmetrische Kurzschlussstrom beträgt näherungsweise
\[
  I_k''=\frac{250\,\mathrm{MVA}}
  {\sqrt{3}\cdot20\,\mathrm{kV}}
  \approx7{,}22\,\mathrm{kA}.
\]

Der Betrag der entsprechenden Ersatzimpedanz ist
\[
  |Z_k|\approx\frac{(20\,\mathrm{kV})^2}{250\,\mathrm{MVA}}
  =1{,}6\,\Omega.
\]

Die Impedanz ist damit nicht vollständig beschrieben: Für eine genaue
Phasorrechnung werden auch der Wirk- und Blindanteil benötigt.

\section{Wirk- und Blindanteil}

Die Kurzschlussimpedanz wird als
\[
  Z_k=R_k+jX_k
\]
geschrieben. Ihr Betrag ist
\[
  |Z_k|=\sqrt{R_k^2+X_k^2}.
\]

Das Verhältnis \(X/R\) beeinflusst den zeitlichen Verlauf des Kurzschlussstroms,
insbesondere den Gleichstromanteil und damit die Stoßstrombeanspruchung.

Wird ausschließlich der Betrag der Impedanz verwendet, kann der anfängliche
symmetrische Wechselstromanteil abgeschätzt werden. Für die Berechnung
des Stoßkurzschlussstroms oder der Ausschaltbeanspruchung ist die
Impedanzphase beziehungsweise das Verhältnis \(X/R\) jedoch wichtig.

\section{Beitrag verschiedener Einspeiser}

In einem realen Netz können Kurzschlussströme aus mehreren Quellen stammen:
\begin{itemize}
  \item dem übergeordneten Versorgungsnetz;
  \item Synchronmaschinen;
  \item Asynchronmaschinen, abhängig von Betriebszustand und Fehlerdauer;
  \item leistungselektronischen Einspeisern mit regelungsabhängigem
        Fehlerstromverhalten.
\end{itemize}

Die Beiträge dürfen nicht unabhängig von ihrer Phasenlage und den
Netzimpedanzen als bloße Beträge addiert werden. Außerdem unterscheiden
sich ihre zeitlichen Verläufe. Für die Bemessung von Schaltgeräten müssen
Anfangskurzschlussstrom, Stoßstrom, Ausschaltstrom und thermische
Beanspruchung entsprechend dem Berechnungsziel getrennt betrachtet werden.

Bei umrichterbasierten Quellen kann ein einfaches Modell als ideale
Spannungsquelle hinter einer festen Impedanz unzureichend sein. Das
Fehlerstromverhalten hängt von Strombegrenzung, Regelung und Schutzstrategie ab.

\section{Thermische und dynamische Beanspruchung}

Der Kurzschlussstrom erzeugt Joule-Verluste, die bei ausreichender Fehlerdauer
zu einer erheblichen Erwärmung führen können. Die elektrodynamischen Kräfte
zwischen stromführenden Leitern steigen näherungsweise mit dem Quadrat
des momentanen Stroms.

Für die Auswahl von Betriebsmitteln müssen deshalb mindestens unterschieden
werden:
\begin{itemize}
  \item der anfängliche symmetrische Kurzschlussstrom;
  \item der Stoßkurzschlussstrom;
  \item der während der Ausschaltzeit wirksame Kurzschlussstrom;
  \item die thermische Beanspruchung über die Fehlerdauer;
  \item das Ausschaltvermögen und die Kurzzeitstromfestigkeit des Betriebsmittels.
\end{itemize}

Die konkrete Bemessung richtet sich nach dem jeweiligen Betriebsmittel,
der Netzkonfiguration und den anwendbaren Normen.

\section{Rechenablauf}

Ein nachvollziehbarer Rechenablauf umfasst:
\begin{enumerate}
  \item Fehlerstelle und betrachteten Betriebszustand festlegen.
  \item Nennspannung und erforderliche Kurzschlusskenngröße bestimmen.
  \item Netz-, Transformator- und Leitungsimpedanzen zusammentragen.
  \item Impedanzen auf eine gemeinsame Spannungsebene beziehen.
  \item Wirk- und Blindanteile komplex addieren.
  \item Den symmetrischen Kurzschlussstrom berechnen.
  \item Bei Bedarf Stoßstrom, Ausschaltstrom und thermische Beanspruchung
        nach dem geeigneten Verfahren bestimmen.
  \item Ergebnis auf Einheiten, Plausibilität und Vollständigkeit prüfen.
\end{enumerate}

\section{Grenzen der vereinfachten Berechnung}

Die einfache Gleichung
\[
  I_k''=\frac{U_{LL}}{\sqrt{3}|Z_k|}
\]
ist ein Ersatzmodell. Sie berücksichtigt nicht automatisch alle
Einflussgrößen einer vollständigen normativen Kurzschlussberechnung.

Je nach Netz und Berechnungsziel können erforderlich sein:
\begin{itemize}
  \item Spannungsfaktoren und Betriebsmittelkorrekturen;
  \item maximale und minimale Netzkonfigurationen;
  \item Widerstandsänderung durch Erwärmung;
  \item Motor- und Generatorbeiträge;
  \item Umrichtermodelle;
  \item die zeitliche Entwicklung des Fehlerstroms;
  \item die tatsächlich vorhandene Fehlerimpedanz.
\end{itemize}

Für Schutzkoordination ist außerdem nicht nur der maximale Kurzschlussstrom
relevant. Der minimale Fehlerstrom kann entscheidend dafür sein, ob ein
Schutzgerät innerhalb der erforderlichen Zeit auslöst.

\section{Übungsaufgaben}

\begin{enumerate}
  \item Berechnen Sie den anfänglichen symmetrischen Kurzschlussstrom
        bei \(U_{LL}=10\,\mathrm{kV}\) und \(S_k''=200\,\mathrm{MVA}\).
  \item Eine Netzimpedanz beträgt \(Z_k=0{,}2+j1{,}0\,\Omega\).
        Bestimmen Sie ihren Betrag und das Verhältnis \(X/R\).
  \item Warum reicht der Betrag der Kurzschlussimpedanz nicht aus,
        um alle Kenngrößen des Kurzschlussverlaufs zu bestimmen?
  \item Erklären Sie den Unterschied zwischen Kurzschlussstrom und
        Kurzschlussleistung.
  \item Warum müssen für die Schutzkoordination sowohl maximale als
        auch minimale Fehlerströme betrachtet werden?
\end{enumerate}

%\chapter{Erdschluss in Mittelspannungsnetzen}

\chapter{Erdschlussberechnung und Erdschlussverhalten}
\label{chap:09-erdschluss}

\section{Einordnung}

Ein Erdschluss liegt vor, wenn ein aktiver Leiter über eine Fehlerstelle
mit Erde oder einem geerdeten leitfähigen Teil verbunden wird. Der
Erdschlussstrom hängt nicht nur von der Netzspannung und der Fehlerimpedanz
ab, sondern wesentlich von der Sternpunktbehandlung, den kapazitiven
Strömen, den Nullsystemimpedanzen und den Erdungswegen.

Die Beurteilung eines Erdschlusses muss deshalb zwischen unterschiedlichen
Netzformen unterscheiden. Insbesondere sind isolierte, kompensierte und
niederohmig geerdete Netze nicht mit demselben vereinfachten Modell
zu behandeln.

\section{Grundlegende Netzformen}

\subsection{Isoliertes Netz}

In einem isolierten Netz ist der Netzsternpunkt nicht absichtlich niederohmig
mit Erde verbunden. Ein einpoliger Erdschluss wird daher wesentlich durch
die Kapazitäten der gesunden Phasen gegen Erde und die übrigen
Ableitpfade bestimmt.

Der Erdschlussstrom kann bei ausgedehnten Kabelnetzen trotz fehlender
direkter Sternpunkterdung beträchtlich sein, da die Kabelkapazitäten
einen kapazitiven Rückstrompfad bilden.

\subsection{Kompensiertes Netz}

In einem kompensierten Netz wird der Sternpunkt über eine
Petersen-Spule mit Erde verbunden. Die induktive Stromkomponente der Spule
wird so eingestellt, dass sie den kapazitiven Erdschlussstrom des Netzes
weitgehend kompensiert.

Die Kompensation ist nicht automatisch vollständig. Netzänderungen,
Schaltzustände, Verluste und Verstimmungen beeinflussen den verbleibenden
Erdschlussstrom. Die tatsächliche Auslegung und Abstimmung müssen zum
Netz und zum vorgesehenen Schutzkonzept passen.

\subsection{Niederohmig geerdetes Netz}

In niederohmig geerdeten Netzen ist der Sternpunkt über eine definierte
Impedanz, häufig einen Widerstand, geerdet. Der Erdschlussstrom wird
dadurch gezielt begrenzt und zugleich ein ausreichend deutliches
Schutzsignal ermöglicht.

Der Widerstand beeinflusst die Höhe des Fehlerstroms, die
Sternpunktverschiebung und die Erdungsspannungen. Die Auslegung muss
Fehlerstrom, zulässige Fehlerdauer, Schutzansprechen und thermische
Beanspruchung gemeinsam berücksichtigen.

\section{Einpoliger Erdschluss und symmetrische Komponenten}

Ein einpoliger Erdschluss an Phase \(a\) kann unter üblichen Voraussetzungen
mit den Sequenznetzen beschrieben werden. Für einen Fehler mit Impedanz
\(Z_f\) gelten am Fehlerpunkt
\[
  I_b=0,\qquad I_c=0.
\]

Daraus folgt für die Sequenzströme
\[
  I_0=I_1=I_2.
\]
Der Fehlerstrom beträgt
\[
  I_f=I_a=I_0+I_1+I_2=3I_0.
\]

Die Sequenznetze werden bei diesem Fehler in Reihe geschaltet. Mit der
Mitsystem-Ersatzspannung \(U_{\mathrm{Th}}\) ergibt sich im vereinfachten
Modell
\[
  I_f=
  \frac{3U_{\mathrm{Th}}}
  {Z_1+Z_2+Z_0+3Z_f}.
\]

Diese Gleichung ist nur dann unmittelbar anwendbar, wenn die verwendeten
Sequenzimpedanzen und die Ersatzspannung die konkrete Netzform abbilden.
Insbesondere darf die Nullsystemimpedanz nicht ohne Prüfung durch die
Mitsystemimpedanz ersetzt werden.

\section{Bedeutung der Nullsystemimpedanz}

Die Nullsystemimpedanz hängt unter anderem ab von:
\begin{itemize}
  \item der Sternpunktbehandlung;
  \item Transformator-Schaltgruppen und Wicklungsverschaltungen;
  \item Erdungsimpedanzen und Erdungsleitern;
  \item Kabelschirmen und metallischen Bewehrungen;
  \item Rückstromwegen über Erde und geerdete leitfähige Strukturen;
  \item dem Aufbau der Netzkomponenten.
\end{itemize}

Eine Erdschlussberechnung, die nur die Mitsystemimpedanz berücksichtigt,
kann deshalb erheblich vom tatsächlichen Verhalten abweichen.

\section{Sternpunktverschiebung und gesunde Phasen}

Bei einem einpoligen Erdschluss in einem isolierten oder kompensierten Netz
kann sich das Potential des Netzsternpunkts gegenüber Erde verschieben.
Dadurch verändern sich die Spannungen der Phasen gegen Erde.

In einem idealisierten isolierten Netz kann die Spannung der beiden
gesunden Phasen gegen Erde bei einem metallischen einpoligen Erdschluss
auf ungefähr die verkettete Spannung ansteigen. Die tatsächlichen
Spannungen hängen von Netzkapazitäten, Dämpfung, Fehlerimpedanz und
Transienten ab.

Diese Spannungsverschiebung ist für die Isolationsbeanspruchung der
gesunden Phasen wichtig. Sie ist ein Grund dafür, dass die Netzform und
das Erdschlusskonzept bei der Auswahl der Isolationskoordination
berücksichtigt werden müssen.

\section{Kapazitiver Erdschlussstrom}

Für eine einzelne Kapazität \(C\) gegen Erde gilt bei sinusförmiger Spannung
\[
  I_C=\omega C U.
\]

Für ein symmetrisches Drehstromnetz ist bei der Berechnung des gesamten
Erdschlussstroms sorgfältig zu unterscheiden, welche Phasenkapazitäten
und welche Phasenspannungen in den jeweiligen Strompfad eingehen.

Für ein idealisiertes isoliertes Netz mit gleicher Phasenkapazität \(C\)
gegen Erde kann der kapazitive Erdschlussstrom eines metallischen
einpoligen Erdschlusses näherungsweise als
\[
  I_E\approx3\omega C U_{\mathrm{ph}}
\]
angegeben werden, wobei \(U_{\mathrm{ph}}\) die normale Strangspannung
gegen Erde und \(C\) die Kapazität einer Phase gegen Erde bezeichnet.

Die Gleichung setzt voraus, dass die drei Phasenkapazitäten gleich sind
und dass die betrachteten Rückstrompfade durch das kapazitive Modell
hinreichend beschrieben werden. Bei realen Kabelnetzen können zusätzliche
Kapazitäten, Schirme, Erdungsverbindungen und Verluste relevant sein.

\subsection{Beispiel}

Ein idealisiertes 20-kV-Netz habe eine gesamte wirksame Phasenkapazität
gegen Erde von
\[
  C=1\,\mathrm{\mu F}
\]
je Phase. Bei \(f=50\,\mathrm{Hz}\) ist
\[
  \omega=2\pi f\approx314{,}16\,\mathrm{rad/s}.
\]

Die Strangspannung beträgt
\[
  U_{\mathrm{ph}}=\frac{20\,\mathrm{kV}}{\sqrt{3}}
  \approx11{,}55\,\mathrm{kV}.
\]

Damit ergibt sich im idealisierten Modell
\[
  I_E\approx3\cdot314{,}16\cdot1\,\mathrm{\mu F}
  \cdot11{,}55\,\mathrm{kV}
  \approx10{,}9\,\mathrm A.
\]

Dieser Wert beschreibt nur den kapazitiven Anteil unter den angegebenen
Annahmen. Er ist nicht ohne Weiteres der gesamte Fehlerstrom eines
realen Netzes.

\section{Kompensation durch Petersen-Spule}

Die Petersen-Spule stellt eine induktive Verbindung zwischen Sternpunkt
und Erde her. Ihre Stromkomponente ist näherungsweise
\[
  I_L=\frac{U_{\mathrm{ph}}}{\omega L}
\]
für eine ideale Induktivität \(L\), an der die Spannung \(U_{\mathrm{ph}}\)
anliegt.

Die Abstimmung zielt darauf, den induktiven Spulenstrom an den kapazitiven
Erdschlussstrom des Netzes anzupassen. In einem vereinfachten Modell
kann der Resonanzzustand durch die Gleichheit der Beträge der
kapazitiven und induktiven Anteile beschrieben werden.

Reale Netze enthalten Verluste und nichtideale Elemente. Deshalb bleibt
ein Wirkstromanteil, und eine vollständige Kompensation ist nicht
automatisch gewährleistet. Zusätzlich muss die Regelung auf
Netzumschaltungen und Änderungen der wirksamen Kapazität reagieren können.

\section{Erdschlussimpedanz und Fehlerort}

Der Erdschluss ist nicht notwendigerweise metallisch. Die Fehlerstelle
kann eine erhebliche Impedanz besitzen, beispielsweise durch Lichtbogen,
verschmutzte Isolatoren, feuchte Oberflächen oder unvollständigen Kontakt
mit Erde.

Eine erhöhte Fehlerimpedanz reduziert häufig den Fehlerstrom. Dadurch kann
die sichere Erkennung durch einen Überstromschutz erschwert werden.
Erdschlussschutzkonzepte nutzen deshalb je nach Netzform unter anderem
Residualstrom, Nullspannung, Richtungsinformation oder weitere
netzabhängige Messgrößen.

Der Fehlerstrom allein genügt nicht in jedem Netz, um Fehlerort und
Fehlerrichtung sicher zu bestimmen.

\section{Schutztechnische Erkennung}

Die Erkennung eines Erdschlusses richtet sich nach dem Netzkonzept.
Mögliche Messgrößen sind:
\begin{itemize}
  \item Summenstrom beziehungsweise Residualstrom;
  \item Nullsystemstrom \(I_0\);
  \item Nullsystemspannung \(U_0\);
  \item Richtungsinformation aus Strom und Spannung;
  \item transiente oder wattmetrische Komponenten, abhängig vom Schutzprinzip.
\end{itemize}

Die konkrete Auswahl hängt von der Sternpunktbehandlung, der Fehlerimpedanz,
dem Umfang des Netzes und den verfügbaren Messwandlern ab. Die
Einstellwerte sind anhand der tatsächlichen Netzparameter und der
Schutzkoordination zu bestimmen.

\section{Erdungsspannung und Berührungsspannung}

Bei Erdschlüssen kann Strom in Erdungsanlagen und den umgebenden Boden
fließen. Dadurch entstehen Potentialgradienten. Zu unterscheiden sind:
\begin{itemize}
  \item Erdpotentialanhebung einer Erdungsanlage;
  \item Berührungsspannung zwischen gleichzeitig berührbaren leitfähigen Teilen;
  \item Schrittspannung zwischen zwei Punkten auf der Erdoberfläche;
  \item über leitfähige Verbindungen eingetragene Potentiale.
\end{itemize}

Die Höhe der Erdpotentialanhebung allein reicht nicht aus, um die
Personengefährdung zu beurteilen. Entscheidend sind unter anderem die
Potentialverteilung, Strompfade durch den Körper, Fehlerdauer,
Oberflächenbeschaffenheit und die konkrete Erdungsgeometrie.

Die Bewertung und Auslegung müssen nach den einschlägigen Normen und
den konkreten Anlagenbedingungen erfolgen. Ein vereinfachtes
Ersatzschaltbild ersetzt keine Erdungsberechnung.

\section{Vergleich der Netzformen}

\begin{center}
\begin{tabular}{p{0.22\linewidth}p{0.34\linewidth}p{0.34\linewidth}}
\hline
\textbf{Netzform} & \textbf{Prägendes Verhalten} & \textbf{Wichtige Prüfgröße}\\
\hline
Isoliert &
Erdschlussstrom wird wesentlich durch Netzkapazitäten und Ableitpfade bestimmt &
Kapazitiver Erdschlussstrom und Sternpunktverschiebung\\
\hline
Kompensiert &
Induktiver Spulenstrom kompensiert den kapazitiven Anteil weitgehend &
Verstimmung, Reststrom und Schutzempfindlichkeit\\
\hline
Niederohmig geerdet &
Sternpunkterdung begrenzt und bestimmt den Erdschlussstrom gezielt &
Fehlerstrom, Fehlerdauer und thermische Beanspruchung\\
\hline
\end{tabular}
\end{center}

Die Tabelle beschreibt allgemeine Tendenzen. Sie ersetzt keine konkrete
Berechnung der jeweiligen Netzform.

\section{Rechenablauf für einen einpoligen Erdschluss}

Ein systematischer Ablauf umfasst:
\begin{enumerate}
  \item Netzform und Sternpunktbehandlung identifizieren.
  \item Fehlerstelle, Fehlerphase und Fehlerimpedanz definieren.
  \item Mit-, Gegen- und Nullsystemimpedanzen ermitteln.
  \item Transformator-Schaltgruppen, Kabelschirme und Erdungswege berücksichtigen.
  \item Die Sequenznetze gemäß den Fehlerbedingungen verschalten.
  \item Den Fehlerstrom und die Nullsystemgrößen berechnen.
  \item Spannungsverschiebung und Schutzgrößen bewerten.
  \item Bei Bedarf Erdpotentialanhebung, Berührungs- und Schrittspannungen
        mit einem geeigneten Erdungsmodell untersuchen.
  \item Ergebnisse gegen die tatsächlichen Anlagenparameter und die
        geltenden Normen prüfen.
\end{enumerate}

\section{Grenzen der vereinfachten Rechnung}

Die analytische Erdschlussrechnung ist nur so gut wie das verwendete
Nullsystemmodell. Insbesondere bei ausgedehnten Kabelnetzen, komplexen
Erdungsanlagen, Lichtbogenfehlern oder leistungselektronischen Einspeisern
kann ein einfaches konzentriertes Ersatzmodell unzureichend sein.

Bei transienten Erdschlüssen spielen zusätzlich Schaltvorgänge, Wanderwellen,
nichtlineare Lichtbogenkennlinien und die frequenzabhängigen Eigenschaften
der Netzkomponenten eine Rolle. Solche Vorgänge erfordern gegebenenfalls
eine zeitbereichsorientierte Simulation.

\section{Übungsaufgaben}

\begin{enumerate}
  \item Erklären Sie, weshalb die Nullsystemimpedanz bei einem Erdschluss
        nicht ohne Weiteres aus der Mitsystemimpedanz abgeleitet werden kann.
  \item Berechnen Sie für ein idealisiertes isoliertes 10-kV-Netz mit
        \(C=0{,}8\,\mathrm{\mu F}\) je Phase gegen Erde den kapazitiven
        Erdschlussstrom bei \(50\,\mathrm{Hz}\).
  \item Beschreiben Sie den Unterschied zwischen isoliertem,
        kompensiertem und niederohmig geerdetem Netz.
  \item Warum kann eine hohe Fehlerimpedanz die Erdschlusserkennung erschweren?
  \item Erläutern Sie, weshalb die Erdpotentialanhebung allein keine
        vollständige Bewertung der Berührungsspannung ermöglicht.
  \item Welche Parameter müssen bekannt sein, bevor eine quantitative
        Erdschlussberechnung für ein reales Mittelspannungsnetz durchgeführt
        werden kann?
\end{enumerate}

\part{Erdung und Schutztechnik}
\chapter{Erdung und Potentialausgleich}
\label{chap:10-erdung-und-potentialausgleich}

Dieses Kapitel ist als Bestandteil des modularen Lehrskripts vorgesehen.
Inhalte werden ergänzt, nachdem die fachlichen Grundlagen und die
Quellenlage für dieses Thema systematisch ausgearbeitet wurden.

\chapter{Schutztechnik}
\label{chap:11-schutztechnik}

Dieses Kapitel ist als Bestandteil des modularen Lehrskripts vorgesehen.
Inhalte werden ergänzt, nachdem die fachlichen Grundlagen und die
Quellenlage für dieses Thema systematisch ausgearbeitet wurden.

\chapter{Überspannungsschutz und Isolationskoordination}
\label{chap:12-ueberspannungsschutz}

Dieses Kapitel ist als Bestandteil des modularen Lehrskripts vorgesehen.
Inhalte werden ergänzt, nachdem die fachlichen Grundlagen und die
Quellenlage für dieses Thema systematisch ausgearbeitet wurden.


\part{Prüfungen und Betriebsmittel}
\chapter{Prüfverfahren und Diagnostik}
\label{chap:13-pruefverfahren}

Dieses Kapitel ist als Bestandteil des modularen Lehrskripts vorgesehen.
Inhalte werden ergänzt, nachdem die fachlichen Grundlagen und die
Quellenlage für dieses Thema systematisch ausgearbeitet wurden.

\chapter{Schaltanlagen und Transformatorenschutz}
\label{chap:14-schaltanlagen-und-transformatorenschutz}

Dieses Kapitel ist als Bestandteil des modularen Lehrskripts vorgesehen.
Inhalte werden ergänzt, nachdem die fachlichen Grundlagen und die
Quellenlage für dieses Thema systematisch ausgearbeitet wurden.


\part{Arbeitssicherheit und Betriebsorganisation}
\chapter{Rollen, Zuständigkeiten und Befugnisse}
\label{chap:15-rollen-und-befugnisse}

Dieses Kapitel ist als Bestandteil des modularen Lehrskripts vorgesehen.
Inhalte werden ergänzt, nachdem die fachlichen Grundlagen und die
Quellenlage für dieses Thema systematisch ausgearbeitet wurden.

\chapter{Begehung, Freischaltung und Erden und Kurzschließen}
\label{chap:16-begehung-freischaltung-euk}

Dieses Kapitel ist als Bestandteil des modularen Lehrskripts vorgesehen.
Inhalte werden ergänzt, nachdem die fachlichen Grundlagen und die
Quellenlage für dieses Thema systematisch ausgearbeitet wurden.

\chapter{Arbeitsfreigabe und Arbeitsdurchführung}
\label{chap:17-arbeitsfreigabe-und-arbeitsdurchfuehrung}

Dieses Kapitel ist als Bestandteil des modularen Lehrskripts vorgesehen.
Inhalte werden ergänzt, nachdem die fachlichen Grundlagen und die
Quellenlage für dieses Thema systematisch ausgearbeitet wurden.

\chapter{Rückgabe und Wiederinbetriebnahme}
\label{chap:18-rueckgabe-wiederinbetriebnahme}

Dieses Kapitel ist als Bestandteil des modularen Lehrskripts vorgesehen.
Inhalte werden ergänzt, nachdem die fachlichen Grundlagen und die
Quellenlage für dieses Thema systematisch ausgearbeitet wurden.


\part{Netzplanung und Praxis}
\chapter{Netzdimensionierung und Netzausbau}
\label{chap:19-netzdimensionierung}

Dieses Kapitel ist als Bestandteil des modularen Lehrskripts vorgesehen.
Inhalte werden ergänzt, nachdem die fachlichen Grundlagen und die
Quellenlage für dieses Thema systematisch ausgearbeitet wurden.

\chapter{Fallstudien}
\label{chap:20-fallstudien}

Dieses Kapitel ist als Bestandteil des modularen Lehrskripts vorgesehen.
Inhalte werden ergänzt, nachdem die fachlichen Grundlagen und die
Quellenlage für dieses Thema systematisch ausgearbeitet wurden.


\appendix
\chapter{Formelverzeichnis}
\label{chap:a-formelverzeichnis}

Dieses Kapitel ist als Bestandteil des modularen Lehrskripts vorgesehen.
Inhalte werden ergänzt, nachdem die fachlichen Grundlagen und die
Quellenlage für dieses Thema systematisch ausgearbeitet wurden.

\chapter{Symbole und Einheiten}
\label{chap:b-symbole-einheiten}

Dieses Kapitel ist als Bestandteil des modularen Lehrskripts vorgesehen.
Inhalte werden ergänzt, nachdem die fachlichen Grundlagen und die
Quellenlage für dieses Thema systematisch ausgearbeitet wurden.

\chapter{Normen und Quellen}
\label{chap:c-normen-und-quellen}

Dieses Kapitel ist als Bestandteil des modularen Lehrskripts vorgesehen.
Inhalte werden ergänzt, nachdem die fachlichen Grundlagen und die
Quellenlage für dieses Thema systematisch ausgearbeitet wurden.


\end{document}
